% Integration fragment. Statement numbers FP.* are temporary identifiers.
% All arguments are AI-produced by GPT-6 Astra (OpenAI, ChatGPT Pro mode).
\section*{Five-point frames: corrections and positive results}
Write $G_n=\widehat\Gamma_{0,[n]}$, $P_n=\widehat\Gamma_{0,n}$ and
$\pi_n:G_n\to\mathfrak S_n$. We use the genuine Drinfeld action, with
half twists $\sigma_i$ and the convention that products act right to left.
Set
\[
 a=\sigma_1\sigma_2\sigma_3\sigma_4,\qquad
 b=\sigma_1\sigma_2\sigma_3,\qquad
 d=\sigma_1^2\sigma_2\sigma_3.
\]
These have respective orders $5,4,3$.

\paragraph{FP.1 (Refuted: arbitrary-multiplier pure frames).}
The assertion that every $\Phi=(\lambda,f)\in\widehat{GT}$ admits
$\alpha_t\in P_5$ with $\Phi(t)=\alpha_t t^\lambda\alpha_t^{-1}$, for
$t=a,b,d$, is false in this normalization.
\emph{Proof.} The standard braid action sends every half twist to a conjugate by a
pure element of its $\lambda$-th power. We use throughout the cocycle
convention of LM7\_PUSH E.65.8, in which the first two generators map to
$\sigma_1^\lambda$ and
$f(\sigma_1^2,\sigma_2^2)^{-1}\sigma_2^\lambda
f(\sigma_1^2,\sigma_2^2)$. Since $\lambda$ is odd, Thus
$\pi_5\Phi=\pi_5$. A pure conjugator would imply
$\pi_5(t)=\pi_5(t)^\lambda$. Since the permutations have orders
$5,4,3$, respectively, this requires $\lambda\equiv1$ modulo the
respective order. Complex conjugation $(-1,1)$ belongs to
$G_{\mathbb Q}\subseteq\widehat{GT}$ and fails all three congruences.
The finite permutation checks are in \texttt{check\_five\_points.py}.
This conclusion concerns the stated genuine action; changing an outer
representative requires specifying its conjugator and changes the displayed
permutation equation.

\paragraph{FP.2 (Known, with derivation: the pentagon has two cocycle classes).}
Let $\rho$ act on $P_5$ by conjugation by
$(\sigma_4\sigma_3\sigma_2\sigma_1)^2$ and put
$s=x_{34}x_{51}^{-1}x_{45}x_{12}^{-1}$. Lochak--Schneps (1997),
Theorem 2 and Proposition 3(iii), give
\[
 f(x_{12},x_{23})=
 \begin{cases}
 \rho(k)^{-1}k,&\lambda\equiv\pm1\pmod5,\\
 \rho(k)^{-1}s k,&\lambda\equiv\pm2\pmod5.
 \end{cases}
\]
The two rows are distinct classes in $H^1(\langle\rho\rangle,P_5)$.
Consequently A.3(b), if asserted without a multiplier restriction, is false:
choose a Galois element whose cyclotomic character is $2$ modulo $5$.
Existence follows from surjectivity of the cyclotomic character on
$\operatorname{Gal}(\mathbb Q(\zeta_5)/\mathbb Q)$.
For multiplier one the first row applies. The primitive is unique because
$P_5^{\rho}=1$, as shown next. This primitive is a path coordinate; its
identification with a particular conjugator of $a$ includes the base-path
conventions and is not an equality of undefined coordinates.

\paragraph{FP.3 (Proved: the three pure centralizers are trivial).}
We have
\[
 C_{P_5}(a)=C_{P_5}(b)=C_{P_5}(d)=1.
\]
\emph{Proof.} For $a$, Lochak--Schneps, Proposition 8(iv) and
Proposition 5(c), applied to $P_5\rtimes\langle a\rangle$, imply that
$\langle a\rangle$ is malnormal in that semidirect product; hence its
centralizer intersects $P_5$ trivially. Replacing $a$ by its square makes no
difference because the two elements generate the same order-five group.

For $b,d$, forget a marked point $p$ fixed by the rotation. This is a map
on the subgroup fixing $p$, not on all of $G_5$. The resulting profinite
Birman sequence has kernel
\[
 1\longrightarrow\widehat F_3\longrightarrow
 (G_5)_p\longrightarrow G_4\longrightarrow1.
\]
The pure version is the completion of the familiar split sequence
$1\to F_3\to\Gamma_{0,5}\to\Gamma_{0,4}\to1$.
Exactness follows directly from characteristic finite-index subgroups of the
finitely generated free kernel: they are invariant under the action of the
free quotient, and the induced finite action factors through a finite
quotient. Passing to the finite extension fixing $p$ retains the same
induced topology on the kernel.

The group $G_4$ has the finite normal Klein group $V_4$ and quotient
$\widehat{\operatorname{PSL}_2(\mathbb Z)}=C_2\amalg C_3$.
The pure subgroup maps injectively to the level-two subgroup of this
quotient. The images of $b,d$ in it have orders $2,3$, respectively.
Centralizers of torsion in this free profinite product meet that level-two
free group trivially. Hence an element of $C_{P_5}(b)$ or $C_{P_5}(d)$ lies
in the kernel $\widehat F_3$.

For $b$, realize the remaining four punctures as $\mu_4$ on
$\mathbb P^1$, with base point $p=\infty$, and $b:z\mapsto iz$.
The orbifold quotient of $\mathbb P^1\setminus\mu_4$ by $C_4$ has one
puncture and two cone points of order four. Its discrete orbifold group is
$C_4*C_4$. The split finite covering sequence therefore identifies
$\widehat F_3\rtimes C_4$ with $C_4\amalg C_4$; the splitting at
$\infty$ is a free factor. The centralizer of that factor is the factor
itself, and its intersection with the free kernel is trivial.

For $d$, take the remaining punctures to be $\{0\}\cup\mu_3$, the
base point $p=\infty$, and $d:z\mapsto\zeta_3z$. The quotient has two
punctures and one cone point of order three. Its group is $C_3*\mathbb Z$,
so its profinite completion is $C_3\amalg\widehat{\mathbb Z}$.
The same factor-centralizer argument applies. These free-product
centralizer facts are the standard torsion-centralizer theorem for free
profinite products: Herfort--Ribes (1985), Theorem B$'$ in Section 4,
printed p.~159, gives $C_G(a)=C_A(a)$ for $1\ne a\in A$ in
$G=A\amalg B$. The opened primary text states in its introductory
Theorem B that the centralizer ``is contained in A''; the displayed
equality is its strengthened formulation. Their Theorem A$'$ and
Theorem 1 also provide the torsion and finite-subgroup conjugacy used
in FP.5 and FP.26.

\paragraph{FP.4 (Proved: full centralizers).}
The full centralizers are $C_{G_5}(a)=\langle a\rangle$,
$C_{G_5}(b)=\langle b\rangle$ and
$C_{G_5}(d)=\langle d\rangle$.
\emph{Proof.} FP.3 makes each centralizer inject into the centralizer of
its permutation. The latter is $C_5$ for $a$ and $C_4$ for $b$, proving
those two assertions. For $d$ it is $C_3\times C_2$, where the extra factor
swaps the two fixed labels. An extra element would therefore give an
involution in $G_5$ with permutation a single transposition. Collas (2012),
Theorem A / Theorem 3.3, makes every profinite involution geometric.
A nontrivial orientation-preserving involution of the sphere has two fixed
points, so on five marked points it must induce two transpositions and one
fixed label. A single transposition with three fixed labels is impossible.

\paragraph{FP.5 (Proved: geometric torsion of order four in $G_5$).}
Every element of order four in $G_5$ is conjugate to a discrete element.
\emph{Proof.} The pure group $P_5$ is torsion free, from its free-by-free
profinite Birman sequence. Thus the permutation of an order-four element
has order four and fixes a unique label $p$. Forget $p$. The free kernel is
torsion free, so the image still has order four in $G_4$.
Every finite-order element of $G_4$ is geometric: its image in
$C_2\amalg C_3$ is conjugate into a factor, and the preimage of either
factor under $G_4\to C_2\amalg C_3$ is the corresponding finite discrete
extension by $V_4$ (orders eight and twelve). The kernel itself is the
geometric $V_4$.
After a conjugation and lift, our element therefore maps to a geometric
order-four rotation of four points. Its preimage is the semidirect product
$\widehat F_3\rtimes C_4=C_4\amalg C_4$ from FP.3. Every finite subgroup
of this free profinite product is conjugate into a discrete factor. Thus
the original element is conjugate to a discrete element of $\Gamma_{0,[5]}$.
Together with Collas's prime-order theorem this also handles all torsion:
an element of order six would have a permutation of type $(3)(2)$, whose
cube is a single transposition, contradicting the prime-order theorem as
in FP.4. The injection of a finite cyclic subgroup into $\mathfrak S_5$
leaves no other orders.

\paragraph{FP.6 (Proved: corrected frame theorem).}
For the stated genuine action, a pure frame at $t=a,b,d$ exists exactly
when $\lambda\equiv1\pmod{\operatorname{ord}(t)}$; if it exists it is
unique and continuous in $\Phi$. In particular all three frames exist
uniquely on the multiplier-one subgroup of $\widehat{GT}$.
\emph{Proof.} Necessity is FP.1. For $a$, Collas's power-conjugacy theorem
supplies $g\in G_5$ with $\Phi(a)=ga^\lambda g^{-1}$. When
$\lambda\equiv1\pmod5$, the permutation of $g$ centralizes the five-cycle,
so multiplication by a power of $a$ makes $g$ pure.

For $b$, FP.5 gives $\Phi(b)=g b^k g^{-1}$, $k=1$ or $3$.
The unique fixed label is the same, so $g$ belongs to the subgroup
$L=(G_5)_p=\widehat{B_4/Z(B_4)}$. Its abelianization is
$\mathbb Z/12$, with every half twist mapping to $1$ and $b$ mapping to $3$.
The GT action on this abelianization is multiplication by $\lambda$.
Consequently $3k=3\lambda\pmod{12}$, so $k\equiv\lambda\pmod4$.
When $\lambda\equiv1\pmod4$, the permutation of $g$ centralizes the
four-cycle and a power of $b$ removes it.

For $d$, use Collas's theorem again:
$\Phi(d)=g d^\lambda g^{-1}$. Its conjugator can permute the two fixed
labels only by swapping them. If it swaps them, multiply it by the
geometric involution $n:z\mapsto1/z$ on
$\{0,\infty\}\cup\mu_3$, which swaps them and inverts $d$.
Then $gn\in L$ and $\Phi(d)=(gn)d^{-\lambda}(gn)^{-1}$.
But $d$ has abelianization $4$ in $L^{\mathrm{ab}}=\mathbb Z/12$.
The last equation would force $4\lambda=-4\lambda\pmod{12}$, which
contradicts invertibility of $\lambda$ modulo three. Thus $g$ fixes both
poles. If $\lambda\equiv1\pmod3$, its permutation is a power of the
three-cycle; multiplication by a power of $d$ makes it pure.
Uniqueness is FP.3. The orbit map $P_5\to G_5$,
$u\mapsto utu^{-1}$, is a continuous injection from a compact space into a
Hausdorff space and hence a homeomorphism onto its image. This proves
continuity of the unique frame.

\paragraph{FP.7 (Proved: crossed composition law; homomorphisms need more).}
On the subgroup where the relevant pure frame exists, with composition
$\Phi\Psi=\Phi\circ\Psi$, one has
\[
 \alpha_t(\Phi\Psi)=\Phi(\alpha_t(\Psi))\alpha_t(\Phi).
\]
Indeed, apply $\Phi$ to the equality defining $\alpha_t(\Psi)$ and use
uniqueness. Reversing the product gives
$\alpha_t(\Psi\Phi)=\Psi(\alpha_t(\Phi))\alpha_t(\Psi)$.
This is a crossed law, not an assertion that a finite image is a
homomorphism. To obtain an ordinary homomorphism into an abelian target one
must prove both the required invariance of the quotient and triviality of
the induced action after the specified correction. A crossed cocycle from
$C_2$ to $C_3$, with the nontrivial element acting by inversion and cocycle
value $1$, already demonstrates that the crossed law alone is insufficient.

\paragraph{FP.8 (Proved: a two-generator presentation).}
Set $h=b^{-1}a$ and $q_i=a^{i-4}ha^{4-i}$ for $1\leq i\leq4$.
A presentation of the discrete group $\Gamma_{0,[5]}$ uses generators
$a,b$ and the following ten relators:
\begin{align*}
 q_iq_{i+1}q_i&=q_{i+1}q_iq_{i+1}&&(1\leq i\leq3),\\
 [q_i,q_j]&=1&&(1\leq i<j\leq4,\ j-i\geq2),\\
 q_4q_3q_2q_1^2q_2q_3q_4&=1,\\
 (q_1q_2q_3q_4)^5&=1,\\
 a&=q_1q_2q_3q_4,\qquad b=q_1q_2q_3.
\end{align*}
This also gives a profinite presentation on taking the closed normal
closure of the relators. To prove it, start from the Artin presentation
with the sphere and full-twist relators. The braid identity
$a\sigma_i a^{-1}=\sigma_{i+1}$ for $i=1,2,3$ and
$b^{-1}a=\sigma_4$ expresses every old generator as $q_i$.
The two last displayed equalities make the maps in both directions
inverse, so this is an explicit Tietze transformation. Also
$d=q_1^2q_2q_3$.
For multiplier one the exact frame constraints are obtained by replacing
$a,b$ in each of these ten relators by
$\alpha_a a\alpha_a^{-1},\alpha_b b\alpha_b^{-1}$ and adjoining
\[
 q_1(\Phi(a),\Phi(b))^2q_2(\Phi(a),\Phi(b))q_3(\Phi(a),\Phi(b))
 =\alpha_d d\alpha_d^{-1}.
\]
These are genuine group equations. Their reduction to a finite quotient
that reads the CM parities is still required before any parity conclusion
can be drawn.

\section*{The harmonic cover and the four cusps of $X_0(6)$}
\paragraph{FP.9 (Proved: abstract harmonic orbifold cover).}
The map $w=4z(1-z)$ is a degree-two orbifold covering from the sphere with
weights $(3,3,2)$ at $(0,1,\infty)$ to the sphere with weights $(3,2,4)$ at
$(0,1,\infty)$. Above $w=0$ lie $z=0,1$ unramified; above $w=1$ lies
$z=1/2$ with index two; and above infinity lies infinity with index two.
The condition $m_{\rm target}=e\,m_{\rm source}$ holds at every point,
and the orbifold Euler characteristics are $1/6$ and $1/12$.
Thus the abstract covering asserted in Lead 6 is correct.
It must not be identified without a comparison with a map of the
principal modular covers $X(3),X(4)$ over the same $j$-line.
Neither principal subgroup contains the other: $T^3$ belongs to
$\Gamma(3)$ but not $\Gamma(4)$, and $T^4$ belongs to
$\Gamma(4)$ but not $\Gamma(3)$. Both are normal in the modular group,
so conjugating them in that group cannot repair the inclusion.
The subgroup $A_4\subset S_4$ gives the abstract spherical-orbifold map;
that fact alone does not compare the two congruence representations or
prove equality of the CM parities.

\paragraph{FP.10 (Known and explicitly checked: $X_0(6)$ and $V_4$).}
A rational Hauptmodul is
\[
 h(\tau)=\frac{\eta(\tau)^5\eta(3\tau)}
 {\eta(2\tau)\eta(6\tau)^5},\qquad u=h+9.
\]
Its four cusp values in the $u$ coordinate are $0,1,9,\infty$.
Thus $Y_0(6)\simeq\mathbb P^1\setminus\{0,1,9,\infty\}$ over
$\mathbb Q$. The Atkin--Lehner involutions are
\[
 w_2(u)=\frac{u-9}{u-1},\qquad
 w_3(u)=\frac9u,\qquad
 w_6(u)=\frac{9(u-1)}{u-9}.
\]
They square to the identity, $w_2w_3=w_3w_2=w_6$, and act on the ordered
cusp list $(0,1,9,\infty)$ as
\[
 (0\ 9)(1\ \infty),\quad (0\ \infty)(1\ 9),\quad
 (0\ 1)(9\ \infty),
\]
respectively. All these identities are checked exactly in the program.
For the modular identification, Maier (2009), Table 2, gives the four
cusp values $0,-8,-9,\infty$ for its reciprocal-scaled Hauptmodul;
Fricke exchanges that coordinate with $h$. Col\`o (2022), Appendix A,
printed p.~245, gives the displayed $h$ and all three involutions directly.
Applying $u=h+9$ gives the formulas above.

\paragraph{FP.11 (Proved correction; comparison remains open).}
The geometric fibre of
$\mathcal M_{0,5}\to\mathcal M_{0,4}$ over the ordered point $9$ is
$\mathbb P^1\setminus\{0,1,9,\infty\}$ and its completed fundamental
group is $\widehat F_3$. This yields the abstract identification with the
completion of $\Gamma_0(6)/\{\pm I\}$. Both descriptions are over
$\mathbb Q$, so the arithmetic fundamental-group identifications respect
the Galois outer action after the chosen geometric base points and paths
are transported. There is no forgetting map from the full unordered
five-point moduli stack to the full unordered four-point stack without
first choosing the point to forget. The group map is on a point stabilizer,
or on the pure group.

The pentagon gives the GT action on the genus-zero tower and compatibility
with forgetting. To identify its action on this fibre with its restriction
from the modular group, one must additionally compare the path or
arithmetic section at $9$. An abstract isomorphism of the free fibre groups
and the identities among the three rational involutions do not establish
that comparison. The unresolved equality is equality, in
$\operatorname{Out}(\widehat F_3)$, of the modular restriction and the
transported fibre action for every $\Phi\in\widehat{GT}$, with the
transport path/section specified. No vanishing of $\Delta$ follows from
FP.10 alone.

\section*{Six and seven points, and the twelve pentagonal structures}
\paragraph{FP.12 (Refuted: universal rigidity in six or seven points).}
Pure centralizers of torsion need not be trivial for six or seven marked
points. On the sphere marked at $\{\pm1,\pm2,\pm3\}$, use the involution
$z\mapsto-z$. On the sphere marked at
$\{0,\pm1,\pm2,\pm3\}$, use the same involution. In either case the
circle $|z|=5/2$ is invariant and essential: it separates four from two
punctures in the first case and five from two in the second.
Its Dehn twist is pure, has infinite order, and commutes with the
involution, because conjugation sends a Dehn twist to the twist about the
image curve. Residual finiteness injects the discrete mapping class group
into its profinite completion, so this remains a nontrivial pure
centralizer there. A frame, when it exists, is therefore a torsor under
this centralizer and is not unique.

\paragraph{FP.13 (Proved: the count twelve).}
A pentagonal structure on a labelled five-element set is a cyclic order
modulo reversal. There are $5!/10=12$, since its stabilizer is the dihedral
group of order ten. There are six subgroups of order five in
$\mathfrak S_5$; each supports two pentagonal structures, represented by
step one and step two around the same five-cycle. Thus twelve counts
unoriented cyclic orders, not order-five cyclic subgroups. This is exactly
the interpretation of cote 148, manuscript folio 29, scan 31, in the
supplied first-pass transcription.
For the geometric configuration $\{\zeta_5^i\}_{i\in\mathbb Z/5}$,
the arithmetic permutation of the five labels is
$i\mapsto\chi(\sigma)i\pmod5$. Its action on the twelve structures
is the ordinary $\mathfrak S_5$ action on
$\mathfrak S_5/D_5$. It has cycle types
\[
 \begin{array}{c|ccc}
 \chi\bmod5&1&2,3&4\\\hline
 \text{cycle type on twelve structures}&1^{12}&2^2 4^2&1^4 2^4.
 \end{array}
\]
In particular the label permutation $i\mapsto-i$ does not fix all
twelve structures.

There is a different, $\mathfrak S_5$-equivariant operation on this
set: replace a generating five-cycle by its $a$-th power, retaining
its cyclic subgroup. Modulo reversal, this operation fixes all twelve
structures for $a=\pm1$ and interchanges the two structures on each
of the six cyclic subgroups for $a=\pm2$, giving six transpositions.
It is not the same as applying the label permutation $i\mapsto ai$.
The exact program enumerates both actions separately and gives their
different cycle types. The former assertion that these actions coincide
has therefore been corrected.

Lochak--Schneps's two-case rule concerns the two special endpoints
corresponding to the two cocycle classes of FP.2, and has the same
$\{\pm1\}$ versus $\{\pm2\}$ dichotomy as the generator-power
operation. That does not identify it with the arithmetic label action
on all twelve cyclic orders. The primitive $k$ describes a path to
the chosen special endpoint; an action on that path space, an action
on the two endpoints, and an action on twelve labelled structures
require separate comparison maps. No such identification is inferred
from their counts alone.

\paragraph{FP.14 (Proved: nontrivial automorphism groups of five-point sets).}
For an unordered five-point subset of $\mathbb P^1(\mathbb C)$, its
full nontrivial group of holomorphic automorphisms is one of
$C_2,C_4,D_3,D_5$. To see this, finite subgroups of
$\operatorname{PGL}_2(\mathbb C)$ are cyclic, dihedral, $A_4,S_4,A_5$.
The latter three have no invariant five-element set: their exceptional
orbit sizes are respectively $(4,4,6)$, $(6,8,12)$ and $(12,20,30)$, and
larger orbits cannot help. For a cyclic rotation of order $m$, nonfixed
orbits have size $m$ and there are at most two fixed points; hence
$5=km+r$, $0\leq r\leq2$. For $m>2$ this gives $m=3,4,5$.
The cases $m=3$ and $5$ automatically admit inversion after normalizing
the single nonfixed orbit to $\mu_m$, giving $D_3,D_5$; the case $m=4$
contains one of the two poles and has group $C_4$. For dihedral groups,
exceptional orbit sizes are $2,m,m$, excluding every other case.
The remaining one-dimensional locus has generic group $C_2$.
This classifies smooth configurations with nontrivial symmetry; the
additional vertices of the 542-vertex manuscript diagram include specified
real, boundary or decorated configurations and are not further isolated
smooth automorphism types.

\paragraph{FP.15 (Open: the 1620 geometric arrows).}
The supplied transcription gives nine orbit sizes summing to $542$ and
$1620$ arrows. If those arrows are taken as independent oriented edge
generators of a connected graph, its free fundamental group has rank
$1620-542+1=1079$. This is merely the graph calculation, not a presentation
of the geometric fundamental groupoid. Such a presentation additionally
requires the actual homotopy class of each arrow and a proof of the
relations among them. The transcription has explicitly uncertain labels
and the scans are not supplied. No claimed fundamental geometric system
of relations is inferred from these counts.

\paragraph{FP.16 (Proved correction: the genus-two fibre needs another branch point).}
The hyperelliptic curve $y^2=x^5-1$ is branched over
$\mu_5\cup\{\infty\}$, six points, and has genus two by
Riemann--Hurwitz. A sphere with the five marked points $\mu_5$ does not by
itself specify this double cover: a hyperelliptic double cover has an even
number of branch points. Choosing a sixth point, or one of the two fixed
poles of the order-five rotation, is additional data. Therefore the
proposed coordinate in a quotient of $\mathbb Z[\zeta_5]$ and the
$\operatorname{Sp}_4(\mathbb Z/12)$ comparison require a proved
functorial extension to that data; they cannot be read from the bare
five-point notation alone.

\paragraph{FP.17 (Proved: the level-six Atkin--Lehner defect is well typed).}
Let $H$ be the closure in $\widehat B$ of
$\Gamma_0(6)/\{\pm I\}$. Subject to the supplied R3 and the standard
arithmetic action on the modular cover $Y_0(6)$, every
$\Phi\in\widehat{GT}_0$ preserves the conjugacy class of $H$.
For multiplier one this follows by reducing the covariance in R3 modulo
six. For arbitrary multiplier choose a Galois element with that multiplier
and write $\Phi$ as a multiplier-one element times it; the Galois element
preserves the subgroup's conjugacy class since $Y_0(6)$ and its map to the
modular orbifold are defined over $\mathbb Q$.

The subgroup $H$ is self-normalizing. Indeed,
$\operatorname{PSL}_2(\mathbb Z/6)\simeq S_3\times A_4$, and the image
of $H$ is $C_2\times C_3$, the product of upper-triangular subgroups.
In $S_3$ the order-two subgroup is self-normalizing; in $A_4$ the
order-three subgroup is self-normalizing. Pullback through the finite
quotient proves the claim. The exact finite enumeration is recorded in
\texttt{check\_normalizers.py}.
Therefore correcting a representative of the outer action to preserve
$H$ produces a well-defined continuous homomorphism
\[
 \rho_6:\widehat{GT}_0\longrightarrow\operatorname{Out}(H).
\]
Two corrections differ by an element of $N_{\widehat B}(H)=H$, so give
the same outer automorphism of $H$. Composition proves that this is a
homomorphism.
For the rational involution $w_e$ of $Y_0(6)$, put
\[
 D_e(\Phi)=\rho_6(\Phi)w_e\rho_6(\Phi)^{-1}w_e^{-1}
 \quad\text{in }\operatorname{Out}(H),\qquad e=2,3,6.
\]
This is a continuous map, trivial on $G_{\mathbb Q}$ because $w_e$ is
rational. Its composition law is
\[
 D_e(\Phi\Psi)=
 \rho_6(\Phi)D_e(\Psi)\rho_6(\Phi)^{-1}D_e(\Phi).
\]
In particular it is a crossed defect, and calling it a finite-order
abelian character requires an additional theorem. Its trivial locus is
the inverse image of the centralizer of $w_e$. Neither
$D_3=1\Longleftrightarrow\Delta=0$ nor $D_6=1$ for all
$\widehat{GT}_0$ is proved by this typing and well-definedness argument.

\paragraph{FP.18 (Proved: cusps alone cannot detect the proposed defect).}
Every $\rho_6(\Phi)$ fixes each of the four cusps of $Y_0(6)$.
The cusp width associated to a double coset
$H g\overline{\langle\tau\rangle}$ is
\[
 [\overline{\langle\tau\rangle}:
  g^{-1}Hg\cap\overline{\langle\tau\rangle}].
\]
It is unchanged by an automorphism preserving the conjugacy class of $H$
and the conjugacy class of the procyclic cusp group, with a unit power on
that group. The four widths are $1,2,3,6$, all distinct; hence no two
cusps can be exchanged. As a result every $D_e(\Phi)$ induces the identity
permutation of the four cusps. Even if the element of R2 has
$D_3\ne1$, that nontriviality must be checked on the free fundamental
group or a stronger quotient; the permutation of the four cusps cannot
witness it.

\paragraph{FP.19 (Proved, matrix statement only: an order-five norm-one target).}
Let $R_N=(\mathbb Z/N)[z]/(z^4+z^3+z^2+z+1)$, with involution
$\bar z=z^{-1}$, and put
\[
 A=\begin{pmatrix}0&0&0&-1\\1&0&0&-1\\0&1&0&-1\\0&0&1&-1\end{pmatrix},
 \qquad
 J=\begin{pmatrix}0&1&0&0\\-1&0&1&0\\0&-1&0&1\\0&0&-1&0\end{pmatrix}.
\]
Then $A^5=I$, $\det J=1$ and $A^{\mathsf T}JA=J$ over the integers.
For every $N$ its centralizer in the group preserving $J$ is canonically
\[
 T_N=\{u\in R_N^{\times}:u\bar u=1\}.
\]
Indeed $e_1,Ae_1,A^2e_1,A^3e_1$ is a basis over $\mathbb Z/N$.
A matrix commuting with $A$ is therefore uniquely a polynomial in $A$ of
degree at most three. Its adjoint for $J$ is obtained by
$A\mapsto A^{-1}$, so the symplectic equation is exactly $u\bar u=1$;
this equation already gives the inverse and no separate unit test is
needed. This argument also works when $N$ is divisible by five.

At $N=12$, exact enumeration of $12^4$ polynomial residues gives
$|T_{12}|=200$, $|T_{12}[2]|=8$, and
$|T_{12}/\{\pm1\}|=100$ with two-torsion of order four. One can also
read these orders from CRT: $R_3=\mathbb F_{81}$ and its norm-one group
has order ten; $R_4$ is the unramified degree-four extension of
$\mathbb Z/4$, whose norm-one group has order twenty. The latter has
principal-unit kernel $\mathbb F_4$ under reduction modulo two.
The four scalar elements $1,5,7,11$ all occur in $T_{12}$.
This proves the advertised form of a possible finite linear target,
but supplies no functorial homomorphism from the five-point mapping class
group to it and no comparison with the actual frame. In particular it
cannot turn a bare five-point primitive into a homomorphism to $T_{12}$.
The enumeration is \texttt{check\_order5\_torus.py} and its output.

\paragraph{FP.20 (Proved: the exact forgotten-point frame map).}
Use the genuine braid action in LM7\_PUSH E.65.8's convention. Let
$L=(G_5)_5$, and let
\[
 q:L\longrightarrow G_4\longrightarrow\widehat B
\]
forget label five and then divide by the geometric Klein four-group.
Writing $u,v$ for the standard braid generators of $\widehat B_3$ modulo
its central full twist, this map is given by
\[
 q(\sigma_1)=u,\qquad q(\sigma_2)=v,\qquad q(\sigma_3)=u.
\]
It follows that
\[
 q(b)=uvu=s,\qquad q(d)=u^2vu=u s=r^{-1},\qquad
 r=uv,\quad sr=u^{-1}.
\]
In particular $x=(sr)^{-2}=u^2$ and $y=sxs=v^2$, exactly the positive
free generators of R3 and R4. The restriction $q:P_5\to F$ is a
continuous surjection, since the pure elements $\sigma_1^2,\sigma_2^2$
map to $x,y$.

The GT tower action respects forgetting. Its action on the quotient
$\widehat B$ is determined by $u,v$ and is
\[
 u\longmapsto u^\lambda,\qquad
 v\longmapsto f(u^2,v^2)^{-1}v^\lambda f(u^2,v^2).
\]
This is the precise action used by E.65.8; at multiplier one it fixes
$sr=u^{-1}$ literally. The Klein subgroup is preserved: the GT action
preserves the permutation map; a torsion element of order two with
permutation a double transposition has trivial image in
$\widehat B=C_2\amalg C_3$, since any nontrivial order-two element there
has nontrivial image in $S_3$. Thus its image remains in $V_4$.
Applying $q$ to the unique frame equations and using the uniqueness of
pure $s,r$ frames in $\widehat B$ gives the exact identities
\[
 \boxed{\ \beta_\Phi=q(\alpha_b(\Phi)),\qquad
 A_\Phi=q(\alpha_d(\Phi))\ }
 \quad(\lambda(\Phi)=1).
\]
The second identity holds because a conjugator of $r^{-1}$ is the same
conjugator of $r$. Hence the CM parities are genuinely read from the
five-point frames after the already specified shear:
\[
 W_\Phi=q_{12}(q(\alpha_b(\Phi)))^{-1}U^{8\rho_\Phi},\qquad
 T_\Phi=q_{12}(q(\alpha_d(\Phi)))^{-1}U^{8\rho_\Phi},
\]
where $(q(\alpha_b(\Phi))^{-1})^{\mathrm{ab}}
=(\rho_\Phi,\rho_\Phi)$. No section at $9$ is needed for these
forgetful identities. The section-comparison gap in FP.11 belongs to
the Atkin--Lehner realization of the fibre.

\paragraph{FP.21 (Refuted: descent of the usual genus-two linear representation).}
The usual genus-two symplectic representation of the five-strand braid
group does not descend to $G_5$ modulo twelve, even after quotienting by
the global scalar sign. Let $J_0$ be the block-diagonal alternating form
with two blocks $\left(\begin{smallmatrix}0&1\\-1&0\end{smallmatrix}\right)$.
For
\[
 v_1=(1,0,0,0)^t,\quad v_2=(0,1,0,0)^t,\quad
 v_3=(1,0,1,0)^t,\quad v_4=(0,0,0,1)^t
\]
put $T_i=I+v_i v_i^tJ_0$. Adjacent transvections satisfy the braid
relation and nonadjacent ones commute, so $\sigma_i\mapsto T_i$
defines a representation of $B_5$ in $\operatorname{Sp}_4(\mathbb Z)$.
The full twist maps to $-I$, but the sphere relator has image
\[
 T_4T_3T_2T_1^2T_2T_3T_4
 =\begin{pmatrix}
 -1&0&0&0\\0&-1&0&0\\0&0&-1&2\\0&0&0&-1
 \end{pmatrix}.
\]
This is not scalar modulo three or twelve. Thus it is not the identity
in either the quotient by $\{\pm I\}$ or by the full scalar center.
The relation required for the five-point sphere is exactly where this
candidate linear comparison fails. It is a representation of the
appropriate six-point hyperelliptic setting with one distinguished point,
and forgetting that point cannot be done silently. This counterexample
refutes this particular proposed genus-two comparison, not the existence
of every conceivable representation of $G_5$ in a symplectic group.

\paragraph{FP.22 (Proved: finite frame ambiguity in the natural five-block target).}
The signed homomorphism $L\to\operatorname{SL}_2(\mathbb Z)$ sending
$\sigma_1,\sigma_3$ to $U$ and $\sigma_2$ to $L_0=
\left(\begin{smallmatrix}1&0\\-1&1\end{smallmatrix}\right)$ is well defined,
since the central full twist of $B_4$ maps to $(UL_0U)^4=I$.
Conjugate the point stabilizer to obtain the analogous representation for
each of the five forgotten labels. On $P_5$, reduction modulo three gives
five homomorphisms $\rho_i:P_5\to\operatorname{SL}_2(\mathbb F_3)$.
Each is onto. Let $H$ be their joint finite image. Then
\[
 H=\{(g_i)\in\operatorname{SL}_2(\mathbb F_3)^5:
          \textstyle\sum_i\epsilon(g_i)=0\},\qquad
 H'=Q_8^5,
\]
where $\epsilon:\operatorname{SL}_2(\mathbb F_3)\to C_3$ sends
$U^2$ to $1$ and has quaternion kernel. In particular
$|H|=8^5\cdot3^4=2654208$.

Here is a proof that avoids enumerating pentagon shadows. For a positive
pure half-twist square $x_{jk}$, the $i$th coordinate of its abelian image
is zero when $i\in\{j,k\}$ and one otherwise. The former follows because
forgetting one of the colliding labels makes its projective image trivial,
so its signed image is central. The latter follows because the twist is
conjugate inside the $i$th stabilizer to $\sigma_1^2$.
The ten vectors with three ones span exactly
$C=\{(c_i)\in\mathbb F_3^5:\sum c_i=0\}$: their pairwise differences
include every $e_i-e_j$.

Modulo the five scalar centers, each quaternion factor becomes the
irreducible two-dimensional $\mathbb F_2$ module $V_i$ for its coordinate
$C_3$. The characters $c\mapsto c_i$ on $C$ are pairwise distinct even
up to inversion, so these five modules are pairwise nonisomorphic.
The image of $H'$ in $\bigoplus V_i$ is a $C$-submodule and projects onto
every $V_i$, since each projection of $H$ is the full
$\operatorname{SL}_2(\mathbb F_3)$. Maschke's theorem applies because
$|C|$ is odd. Multiplicity one then makes that image all of
$\bigoplus V_i$. Choose two elements of $H'$ whose images are a basis
of a single $V_i$ and zero in the other four factors. Their commutator is
$-I$ in factor $i$ and $I$ elsewhere. Thus $H'$ contains all five scalar
centers as well, proving $H'=Q_8^5$ and the asserted formula for $H$.

Induce the signed stabilizer representation to the full group $G_5$.
This gives a genuine finite five-block representation modulo twelve.
An element of $P_5$ acts block diagonally, and $b$ fixes the block with
forgotten label five, acting there by $S=UL_0U$.
The image of $P_5'$ is trivial modulo four in every block: a pure element
maps into the kernel of
$\operatorname{SL}_2(\mathbb Z/4)\to\operatorname{SL}_2(\mathbb F_2)$,
which is abelian. Combining this with $H'=Q_8^5$ and CRT shows that the
image of $P_5'$ contains the element supported only in the fifth block
with value
\[
 \widetilde S=\begin{pmatrix}9&4\\8&9\end{pmatrix}\pmod{12}.
\]
It is $I$ modulo four and $S$ modulo three, and it commutes with the
finite image of $b$. Multiplying a finite $b$-frame by this centralizer
element changes $\epsilon_2(W)$ by one without changing the frame
conjugacy equation. Similarly the element supported in the fifth block
with value $5I$ centralizes the finite image of $d$ and changes
$\epsilon_3(T)$ by one. Both toggles have lifts in $P_5'$, so their
forgotten-point abelian coordinates and hence the shear coordinate are
zero.

Thus even the identity automorphism of this quotient admits solutions of
the finite frame equations with either apparent CM parity. The
presentation equations of FP.8 cannot disambiguate these finite primitive
choices, because the images of $a,b,d$ under the represented automorphism
have not changed. This is a precise failure of this proposed finite
reading of frames. It is not a mixed-sign element of $\widehat{GT}$, not
an assertion about extendability of the finite primitive choices, and not
a disproof that $\Delta$ factors through some action quotient with the
correct extendability condition. A finer finite quotient in which the
images of actual global primitives are separated is still needed.
The exact small group and module checks, including the two CRT matrices,
are recorded in \texttt{check\_five\_block\_modules.py} and its output.

\paragraph{FP.23 (Proved: a constructive finite quotient that separates frames).}
Let $B^\wedge=C_2\amalg C_3$, with generators $s,r$, and let
$q:B^\wedge\twoheadrightarrow Q$ be a continuous finite quotient in which
$s,r$ retain their respective orders. Write $H_2=\langle q(s)\rangle$,
$H_3=\langle q(r)\rangle$, and form the finite left permutation modules
\[
 M_2=\mathbb F_2[Q/H_2],\qquad M_3=\mathbb F_3[Q/H_3].
\]
Put $E=(M_2\oplus M_3)\rtimes Q$, with its ordinary semidirect-product
multiplication. If $e_2,e_3$ denote the basis vectors at the identity
cosets, define
\[
 S=(e_2,0,q(s)),\qquad R=(0,e_3,q(r)).
\]
Since the relevant cyclic subgroup fixes its identity coset, $S^2=R^3=1$.
The universal property of the free profinite product gives a continuous
homomorphism $\theta:B^\wedge\to E$ lifting $q$. Its finite image
$E_0=\langle S,R\rangle$ satisfies
\[
 \operatorname{pr}_Q C_{E_0}(S)=H_2,\qquad
 \operatorname{pr}_Q C_{E_0}(R)=H_3.
\]

We prove the first identity; the second is identical with three in place
of two. If $(m_2,m_3,g)$ centralizes $S$, then $g$ centralizes $q(s)$ and
\[
 (1-q(s))m_2=(1-g)e_2.
\]
Because $g$ centralizes $q(s)$, both cosets $H_2$ and $gH_2$ are fixed by
$q(s)$. In the coinvariants $M_2/(1-q(s))M_2$, basis vectors are indexed
by orbits of $\langle q(s)\rangle$ on $Q/H_2$. These two fixed cosets
therefore give distinct, independent basis vectors unless $gH_2=H_2$.
The displayed equation forces $g\in H_2$. Conversely $S^k$ centralizes
$S$ and projects to every element of $H_2$.

Take now $Q=\operatorname{SL}_2(\mathbb Z/12)/\{\pm I\}$, of order
$576$, and let $F$ be the level-two subgroup. If two elements of $E$ with
$Q$-projections in $q(F)$ conjugate $S$ to the same element, their
$Q$-projections are equal. Indeed their ratio centralizes $S$, its
projection lies in $H_2$, and
$H_2\cap q(F)=1$, since the map $Q\to S_3$ is injective on $H_2$.
The analogous assertion holds for $R$. Thus these two lifted finite
conjugacy equations determine the full $Q$-residues of the actual pure
frames, not merely their conjugacy classes in $Q$.

This construction has ambient order
$576\cdot2^{288}\cdot3^{192}$. It is not claimed minimal.
The exact program checks all $576$ elements of $Q$, the cyclic cosets,
and the fixed-coset obstructions. In $Q$ the centralizers of $q(s),q(r)$
have orders $16,18$; the construction excludes respectively $14,15$
unwanted centralizer residues. The large module group need not be
explicitly enumerated to verify this proof.

\paragraph{FP.24 (Proved: characteristic detector and evaluation on the multiplier-one group).}
There is an explicitly defined finite characteristic quotient of
$B^\wedge$ through which the genuine multiplier-one action determines
$\delta$ of R4. One construction is
\[
 K_* =\bigcap_{h\in\operatorname{Hom}_{\rm cont}(B^\wedge,E)}\ker h,
 \qquad Q_*=B^\wedge/K_*.
\]
The indexing set is finite: its elements are exactly pairs of elements
of $E$ whose squares and cubes, respectively, are one. Thus $K_*$ is
open, normal and characteristic; precomposition by an automorphism
permutes the finite set of homomorphisms. It lies in $\ker\theta$,
so $Q_*$ maps to $E_0$. Equivalently one may use the characteristic core
of $\ker\theta$, which is open because a finitely generated profinite
group has only finitely many subgroups of a specified finite index.

We use the \emph{genuine} action selected by E.65.8,
$u\mapsto u^\lambda$ and
$v\mapsto f(u^2,v^2)^{-1}v^\lambda f(u^2,v^2)$, with
$s=uvu$, $r=uv$, $sr=u^{-1}$.
This is an explicit representative, not an unspecified outer class.
Let $\rho:\widehat{GT}_0\to\operatorname{Aut}(Q_*)$ be the resulting
continuous homomorphism. For $\Phi\in\mathcal H$, project
$\rho(\Phi)(s),\rho(\Phi)(r)$ into $E_0$. Solve the two finite
conjugacy equations there, restricting the conjugators to have
$Q$-projection in $q(F)$. Actual pure frames guarantee existence.
FP.23 guarantees that the two $Q$-projections obtained are unique and
are $q(\beta_\Phi),q(A_\Phi)$.

Their signed lifts in the specified image of $F$ in
$\operatorname{SL}_2(\mathbb Z/12)$ are unique: modulo four that image
is the displayed $V_4$ of R3, which does not contain $-I$.
For the actual multiplier-one frames the lift is $I$ modulo four.
Write $\epsilon:\operatorname{SL}_2(\mathbb F_3)\to\mathbb F_3$
for the abelianization with $\epsilon(U^2)=\epsilon(L_0^2)=1$.
Since $\beta_\Phi^{\rm ab}=(-\rho_\Phi,-\rho_\Phi)$,
\[
 \epsilon(q_3(\beta_\Phi))=-2\rho_\Phi=\rho_\Phi\pmod3.
\]
Consequently its residue already determines the shear $U^{8\rho_\Phi}$
modulo twelve. We can compute $T,W$ and their two parities by R3--R4.
This defines a well-defined character
$d:\rho(\mathcal H)\to\mathbb F_2$ with
$\delta=d\circ\rho|_{\mathcal H}$.
It is not asserted that every automorphism of $Q_*$ lies in
$\rho(\mathcal H)$; the detector evaluates actual action residues,
and finite admissibility is not a lifting theorem.

\paragraph{FP.25 (Proved: reduction for arbitrary multipliers; arithmetic table not computed).}
Retain $\rho$ and $d$ from FP.24, and put
\[
 J=\chi\bigl(G_{\mathbb Q}\cap\ker\rho\bigr)
       \subset\widehat{\mathbb Z}^{\times}.
\]
This is open: $G_{\mathbb Q}\cap\ker\rho$ is open, and the surjective
continuous cyclotomic homomorphism of profinite groups is open.
Choose a positive integer $M$ such that
$\ker(\widehat{\mathbb Z}^{\times}\to(\mathbb Z/M)^\times)\subset J$.
For each $a\in(\mathbb Z/M)^\times$, choose a Galois element $\sigma_a$
whose multiplier has residue $a$; retain only the finite table
$a\mapsto\rho(\sigma_a)$. Then for every
$\Phi\in\widehat{GT}_0$,
\[
 \Delta(\Phi)=d\bigl(\rho(\Phi)\rho(\sigma_a)^{-1}\bigr),
 \qquad a=\lambda(\Phi)\bmod M.
\]
The argument of $d$ lies in $\rho(\mathcal H)$ even though
$\Phi\sigma_a^{-1}$ need not have multiplier exactly one.
Indeed its multiplier lies in $J$, so choose
$\sigma_0\in G_{\mathbb Q}\cap\ker\rho$ with that exact multiplier.
Then $h=\Phi\sigma_a^{-1}\sigma_0^{-1}$ lies in $\mathcal H$ and
$\rho(h)=\rho(\Phi)\rho(\sigma_a)^{-1}$. R4 gives
$\Delta(\Phi)=\delta(h)$. This also proves independence of the choices.

Adjoin a characteristic finite quotient covering the congruence quotient
modulo $M$. Since $sr$ has order $M$ in that congruence quotient and the
genuine action sends it to its $\lambda$-th power, the enlarged action
quotient records $\lambda\bmod M$ as well. Thus it detects the full
character. This is an exact reduction; neither $M$, the finite Galois
table, the smallest detector, nor all pentagon residues at that detector
have been computed here. Equivalently $M$ may be chosen as a cyclotomic
conductor for the abelian intersection of the finite Galois extension
cut out by $\rho|_{G_{\mathbb Q}}$.

\paragraph{FP.26 (Proved: the multiplier-one quartic frame without E.39).}
The quartic frame needed in LM7\_PUSH E.64.7 exists uniquely on
$\widehat{GT}\cap\mathcal H$. This assertion does not require the
unverified general periodic-braid statement E.39. We now prove it,
including its exact linear comparison, in the genuine normalization of
E.65.8. In this paragraph only, write
\[
 B_4=\langle a,b,c\mid aba=bab,\ bcb=cbc,\ ac=ca\rangle,
 \qquad d=abc,
\]
so that $d$ is the order-four element called $b$ after passage to $G_5$
in FP.1--FP.6. Let
\[
 q:\widehat B_4\longrightarrow\widehat B_3,
 \quad q(a)=q(c)=a,\quad q(b)=b,\qquad
 s=aba,\quad z_3=(ab)^3=s^2,\quad z_4=d^4.
\]
Let $j$ be the section generated by the first three strands. Put
\[
 \xi_1=ca^{-1},\qquad
 \xi_2=aba c^{-1}b^{-1}a^{-1},\qquad K=\ker q.
\]
For every $\Phi\in\widehat{GT}\cap\mathcal H$, there is exactly one
$H_\Phi\in\widehat B_4$ such that
\[
 \Phi_4(d)=H_\Phi dH_\Phi^{-1},\qquad q(H_\Phi)=\beta_\Phi.
 \tag{FP.26.1}
\]
Here $\beta_\Phi$ is the unique harmonic pure frame lifted to the free
complement $\langle a^2,b^2\rangle$ in $\widehat B_3$.

\emph{The fold and its free kernel.}
The following exact Nielsen action is obtained from the three Artin
relations:
\[
 \begin{array}{c|cc}
 &\xi_1&\xi_2\\ \hline
 a&\xi_1&\xi_2\xi_1\\
 b&\xi_1\xi_2^{-1}&\xi_2.
 \end{array}
 \tag{FP.26.2}
\]
These two automorphisms of $F(\xi_1,\xi_2)$ satisfy the braid relation.
Conversely, in the resulting semidirect product $F_2\rtimes B_3$, set
$c=\xi_1a$. The elements $a,b,c$ satisfy the three Artin relations:
$a$ fixes $\xi_1$, and
\[
 bcb=\xi_1\xi_2^{-1}aba=cbc,
\]
because $(ab)\xi_1(ab)^{-1}=\xi_2^{-1}$.
Moreover the displayed braid word defining $\xi_2$ gives exactly that
free generator. Thus the two presentations give mutually inverse
homomorphisms, proving $B_4=F(\xi_1,\xi_2)\rtimes B_3$ and
$K=F(\xi_1,\xi_2)$ without importing a Ferrari-kernel assertion.
Every finite characteristic quotient of this finitely generated free
kernel has a finite automorphism group; its action therefore factors
through a finite quotient of $B_3$. These characteristic quotients are
cofinal, so completion gives the exact split extension
\[
 \widehat B_4=\widehat K\rtimes\widehat B_3.
 \tag{FP.26.3}
\]

\emph{Equivariance in the fixed normalization.}
Write $f_0=f(a^2,b^2)$ and $e=f(z_3,c^2)$. The multiplier-one
Drinfeld formulas on these literal generators are
\[
 \Phi_4(a)=a,\qquad
 \Phi_4(b)=f_0^{-1}bf_0,\qquad
 \Phi_4(c)=e^{-1}ce.
 \tag{FP.26.4}
\]
For the last formula the usual first argument is $ba^2b$.
The identity $z_3=a^2(ba^2b)$ permits its replacement by $z_3$:
$a^2$ commutes with both arguments, and $f$ has zero exponent sum in
each variable. Since $q(z_3)=z_3$ commutes with $a^2$,
$q(e)=f(z_3,a^2)=1$. Consequently
\[
 q\Phi_4=\Phi_3q,\qquad \Phi_4j=j\Phi_3.
 \tag{FP.26.5}
\]
This proves preservation of $\widehat K$ as well. Collas (2012),
Section 1, records the standard braid formulas in his opposite
conjugator convention; our convention here is fixed by (FP.26.4).
His prime-order theorem alone does not supply the quartic claim.
The order-four input is instead the argument of FP.5--FP.6.

\emph{Lifting the five-point frame.}
The quotient $B_4/\langle z_4\rangle$ is the label-five stabilizer in
$\Gamma_{0,[5]}$. Its completion is the corresponding open subgroup
$L$ of $G_5$, and the central kernel is exactly
$\overline{\langle z_4\rangle}\cong\widehat{\mathbb Z}$.
The last isomorphism follows also from exponent sum, which takes
$z_4$ to $12$: the induced topology on this infinite cyclic group is
its full profinite topology. By FP.6 there is a pure element
$\bar H\in P_5\subset L$ conjugating $\bar d$ to
$\Phi(\bar d)$. Lift it to $H_0\in\widehat B_4$. For some
$t\in\widehat{\mathbb Z}$ one has
\[
 \Phi_4(d)=H_0dH_0^{-1}z_4^t.
\]
The exponent-sum homomorphism $\widehat B_4\to\widehat{\mathbb Z}$
is fixed by a multiplier-one action. The two sides have exponent sums
$3$ and $3+12t$, so $12t=0$; multiplication by $12$ is injective on
$\widehat{\mathbb Z}$, hence $t=0$. This proves exact conjugacy
upstairs, including the central power which a projective argument alone
would leave undetermined.

The harmonic frame also lifts exactly to $\widehat B_3$:
its possible central discrepancy is killed by the same exponent-sum
argument, now with $z_3$ of exponent sum six. Applying (FP.26.5) gives
\[
 q(H_0)sq(H_0)^{-1}=\beta_\Phi s\beta_\Phi^{-1}.
\]
The centralizer of $s$ in $\widehat B_3$ is
$\overline{\langle s\rangle}\cong\widehat{\mathbb Z}$.
Indeed its image in $C_2\amalg C_3$ lies in the order-two factor by
the free-profinite-product centralizer theorem; its remaining factor
is a central power of $z_3=s^2$. Therefore
$q(H_0)=\beta_\Phi s^a$ for a unique $a\in\widehat{\mathbb Z}$.
Replacing $H_0$ by $H_0d^{-a}$ proves existence in (FP.26.1).

\emph{Uniqueness and the quartic group.}
Equations (FP.26.2) give
\[
 d=\xi_2^{-1}s,\qquad
 d\xi_1d^{-1}=\xi_2^{-1},\qquad
 d\xi_2d^{-1}=\xi_1.
 \tag{FP.26.6}
\]
The preimage of $\overline{\langle s\rangle}$ is consequently
$\widehat K\rtimes\overline{\langle d\rangle}$, with the displayed
order-four action. If $y=b^{-1}d^2b$, direct braid multiplication gives
$yd^{-2}=\xi_2^{-1}$ and $y^2=d^4$. With
$k_0=\xi_2^{-1}$, $k_1=\xi_1^{-1}$, the action is
$k_0\mapsto k_1$, $k_1\mapsto k_0^{-1}$.
Thus the mutually inverse presentation maps
\[
 \langle D,Y\mid Y^2=D^4\rangle
 \ \rightleftarrows\ F(k_0,k_1)\rtimes\mathbb Z,
 \quad D\longmapsto d,\quad Y\longmapsto k_0d^2
\]
identify this preimage with the profinite completion of the quartic
group. Its quotient by $\overline{\langle d^4\rangle}$ is
$C_4\amalg C_2$. The image of $\widehat K$ is the kernel of the
map to $C_4$ that sends $D$ to $1$ and $Y$ to $2$.
If $k\in\widehat K$ centralizes $d$, its image in this free profinite
product lies in the factor $\langle\bar d\rangle$ and also in that
kernel, so it is trivial. Then $k\in\overline{\langle d^4\rangle}$,
and $q(k)=1$ forces $k=1$, since $q(d^4)=s^4$ has infinite
profinite order. Hence
\[
 C_{\widehat K}(d)=1.
 \tag{FP.26.7}
\]
Two normalized frames have ratio in this centralizer, proving uniqueness.
Alternatively FP.4, followed by the fold, gives the same conclusion.
The normalized frame is continuous: the map
$H\mapsto(HdH^{-1},q(H))$ is a continuous injection from a compact
space onto its image after (FP.26.7).

Uniqueness and (FP.26.5) also give the exact crossed law
\[
 H_{\Phi\Psi}=\Phi_4(H_\Psi)H_\Phi.
\]
Consequently $\operatorname{Inn}(H_\Phi^{-1})\Phi_4$ is an actual
action on the quartic subgroup, fixing $d$ and $z_4$.

\emph{The signed linear comparison.}
On $\widehat K^{\mathrm{ab}}$ in the ordered basis
$(\xi_1,\xi_2)$, (FP.26.2) and (FP.26.6) give
\[
 Q(a)=U=\begin{pmatrix}1&1\\0&1\end{pmatrix},\qquad
 Q(b)=L_0=\begin{pmatrix}1&0\\-1&1\end{pmatrix},\qquad
 Q(d)=S=UL_0U.
\]
Write $M_\Phi$ for the action of $\Phi_4$ on this homology.
Both arguments of $e$ commute with $a$, so (FP.26.4) gives
$\Phi_4(\xi_1)=e^{-1}\xi_1e$ with $e\in\widehat K$.
Thus $M_\Phi e_1=e_1$. Covariance at $a$ gives
$M_\Phi U=UM_\Phi$, which implies
\[
 M_\Phi=\begin{pmatrix}1&t_\Phi\\0&1\end{pmatrix},\qquad
 M_\Phi L_0M_\Phi^{-1}
       =Q(f_0)^{-1}L_0Q(f_0).
\]
Because $q(H_\Phi)=\beta_\Phi$, (FP.26.3) writes it uniquely as
$H_\Phi=j(\beta_\Phi)k_\Phi$ with $k_\Phi\in\widehat K$.
Inner conjugation by $k_\Phi$ acts trivially on
$\widehat K^{\mathrm{ab}}$. The normalized quartic action therefore
has the literal matrix
\[
 \boxed{\ C_2(\Phi)=Q(\beta_\Phi)^{-1}M_\Phi\ },\qquad
 C_2(\Phi)S=SC_2(\Phi).
 \tag{FP.26.8}
\]
No sewing theorem, unspecified path, or assumed quartic instance of E.39
enters this identity. It proves the quartic input used in E.64 at
multiplier one. It does not by itself prove E.39 for arbitrary numbers
of strands or arbitrary multipliers, nor identify the normalized quartic
action with a separate orbifold action of the same GT pair.
The discrete identities and the integral homology signs are independently
verified by the faithful Artin representation in
\texttt{check\_quartic\_frame\_input.py}; its complete output is retained.

\paragraph{FP.27 (Proved: the explicit order-four and order-three cocycles).}
Return to the global notation of FP.1, and write
$u_i=\sigma_i$ for $1\leq i\leq4$ and
$z_i=(u_1\cdots u_{i-1})^i$ for the full twist on the first $i$
strands. For $\Phi=(1,f)\in\widehat{GT}$ put
\[
 F_i=f(z_i,u_i^2)\quad(2\leq i\leq4),\qquad F_1=1.
\]
All these words are read in the specified sphere quotient of the braid
group. The formula with $z_i$ agrees with the usual Drinfeld first
argument: their ratio is $z_{i-1}$, which commutes with that argument
and with $u_i^2$, so it can be inserted in a commutator word $f$.
In the convention of E.65.8,
\[
 \begin{split}
 E_b&=u_1F_2^{-1}u_2F_2F_3^{-1}u_3F_3
                         u_3^{-1}u_2^{-1}u_1^{-1},\\
 E_d&=u_1^2F_2^{-1}u_2F_2F_3^{-1}u_3F_3
                         u_3^{-1}u_2^{-1}u_1^{-2},\\
 E_a&=\left(\prod_{i=1}^{4}F_i^{-1}u_iF_i\right)a^{-1}.
 \end{split}
 \tag{FP.27.1}
\]
These are exactly $E_t=\Phi(t)t^{-1}\in P_5$.
Let $\vartheta_t$ be conjugation by $t$, of order $m=5,4,3$ respectively.
The defining torsion relation gives the nonabelian cocycle condition
\[
 E_t\,\vartheta_t(E_t)\cdots\vartheta_t^{m-1}(E_t)=1.
 \tag{FP.27.2}
\]
The unique pure frames of FP.6 are the unique primitives, in the
explicit convention
\[
 E_t=\alpha_t\vartheta_t(\alpha_t^{-1}).
 \tag{FP.27.3}
\]
Indeed $(E_tt)^m=\Phi(t)^m=1$ proves (FP.27.2), and multiplying the
frame equation by $t^{-1}$ proves (FP.27.3). Conversely (FP.27.3)
recovers that frame equation. This is an explicit answer for the
order-four and order-three primitive cocycles; it does not identify an
arbitrary-multiplier Lochak--Schneps path primitive with an unspecified
torsion frame. The latter assertion was refuted in FP.1--FP.2.

There is also an exact relation between the two local frames which is
already visible before the fifth generator is used. Since $d=u_1b$ and
$\Phi(u_1)=u_1$,
\[
 \alpha_d d\alpha_d^{-1}
       =u_1\alpha_b b\alpha_b^{-1}.
 \tag{FP.27.4}
\]
If $h=\alpha_b^{-1}\alpha_d$, this is equivalently
\[
 b h b^{-1}=\alpha_b^{-1}u_1^{-1}\alpha_b h u_1.
 \tag{FP.27.5}
\]
Under the exact forgetful map of FP.20 it becomes the corresponding
relation between $A,\beta$ in $C_2\amalg C_3$, with $q(b)=s$,
$q(d)=r^{-1}$ and $q(u_1)=(sr)^{-1}$. Every element of
$\mathcal H\subset\widehat{GT}_0$ already satisfies that four-point
relation. The verified R2 mixed-sign element therefore rules out
deducing equality of the two parities from (FP.27.4) or (FP.27.5)
alone. This pinpoints why the remaining relations involving the fifth
point are essential.

\paragraph{FP.28 (Proved: an actual finite five-point level separates the two frame residues).}
The detector of FP.23 can be embedded in one explicitly specified finite
characteristic quotient of $G_5$. This makes precise the finite level
needed for A.3(e), without claiming a parity-vanishing enumeration.
Let $L=(G_5)_5$, let $q:L\to C_2\amalg C_3$ be the forgetful map
of FP.20, and let $\theta:C_2\amalg C_3\to E_0$ be FP.23.
Set
\[
 N=P_5\cap\ker(\theta q),\qquad m=[G_5:N],\qquad
 N_*=\bigcap_{h\in\operatorname{Hom}_{\rm cont}
                         (G_5,\mathfrak S_m)}\ker h.
 \tag{FP.28.1}
\]
The subgroup $N$ is open in $L$, and hence in $G_5$, because $L$
has index five. The final intersection is finite: a homomorphism is
determined by the images of the two generators $a,b$. It is invariant
under every automorphism of $G_5$. The coset action on $G_5/N$ is one
of these homomorphisms, so $N_*$ lies in its kernel, the normal core
of $N$. Consequently $N_*$ is open and characteristic, is contained
in $N$, and
\[
 \mathcal Q_5=G_5/N_*
\]
is a finite characteristic quotient. Since $N_*\subset P_5$, it
retains the permutation map to $\mathfrak S_5$. The image $L_*$ of
$L$ admits a specified homomorphism $L_*\to E_0$, since
$N_*\subset\ker(\theta q)$.

For an actual $\Phi\in\widehat{GT}\cap\mathcal H$, its genuine
action on $\mathcal Q_5$ determines the residues of both frames in
$Q=\operatorname{SL}_2(\mathbb Z/12)/\{\pm I\}$.
To see this, solve
\[
 \Phi(b_*)=h_b b_*h_b^{-1},\qquad
 \Phi(d_*)=h_d d_*h_d^{-1}
 \tag{FP.28.2}
\]
with $h_b,h_d$ in the image of $P_5$ in $\mathcal Q_5$.
Actual frames supply solutions. Two choices differ by a centralizer.
Projecting the equations into $E_0$ gives the conjugacy equations
for $S$ and $R^{-1}$, respectively. The uniqueness result of FP.23
then forces their $Q$-projections to be
$q_{12}(\beta_\Phi)$ and $q_{12}(A_\Phi)$. Their signed lifts and
the shear are recovered exactly as in FP.24. Thus the two individual
parities and $\Delta(\Phi)$ are functions of this genuine finite
five-point action residue.

The construction is finite and effective from the presentations:
enumerate the finite group $E$, the subgroup generated by $S,R$,
the coset action of $G_5$ on $G_5/N$, and then all pairs of
permutations in $\mathfrak S_m$ satisfying the ten relators of FP.8.
The product of their homomorphisms has kernel $N_*$.
It is a very large bound, not the smallest level. The representation
relations of FP.8 can now be tested at a level which really determines
the residues of actual frames; the ambiguity exhibited in FP.22 is
absent after projection to $Q$. No enumeration of all pentagon shadows
in $\mathcal Q_5$ is claimed here. The unresolved step is to show
that all admissible pentagon action residues in such a detector have
equal parities, or to construct a compatible sequence of mixed-sign
residues giving an actual element. An isolated finite mixed-sign
solution would prove neither alternative.

\paragraph{FP.29 (Proved: the explicit detector also removes the outer-action gauge).}
The characteristic quotient $Q_*$ of FP.24 detects $\delta$ from the
\emph{outer} action on the multiplier-one subgroup. The extra fact is
that the same module extension separates the cusp centralizer. Put
$\tau=sr$, $\bar\tau=q(\tau)$ and $T_\tau=SR=\theta(\tau)$.
Then
\[
 \operatorname{pr}_Q C_{E_0}(T_\tau)=\langle\bar\tau\rangle,
 \qquad |\bar\tau|=12.
 \tag{FP.29.1}
\]
To prove this, the $M_2$ component of $T_\tau$ is $e_2$.
If $(m,g)$ centralizes $T_\tau$, its projection $g$ centralizes
$\bar\tau$ and
\[
 (1-\bar\tau)m_2=(1-g)e_2.
\]
Passing to $\langle\bar\tau\rangle$-coinvariants of $M_2$ says that
the cosets $H_2$ and $gH_2$ belong to the same $\bar\tau$-orbit.
Thus $g\in\langle\bar\tau\rangle H_2$. But $H_2$ intersects
$C_Q(\bar\tau)$ trivially: the nonidentity element $q(s)$ does not
commute with $q(sr)$, as its displayed matrices verify. It follows
that $g\in\langle\bar\tau\rangle$. Powers of $T_\tau$ give the
reverse inclusion, proving (FP.29.1).

Now suppose $\Phi\in\mathcal H$ acts by an inner automorphism on
$Q_*$. Choose its finite inner conjugator $g_*$. Since the genuine
action fixes $\tau$ literally, the image of $g_*$ in $E_0$ centralizes
$T_\tau$. Thus its $Q$-projection is $\bar\tau^k$ for some
$k\in\mathbb Z/12$. The genuine action also fixes the permutation
quotient $S_3$ pointwise; hence that projection is pure, since
$Z(S_3)=1$. The two actual pure frames and $g_*$ conjugate $S,R$ in
the same way. FP.23 then gives
\[
 q(\beta_\Phi)=q(A_\Phi)=\bar\tau^k.
\]
In the signed matrix normalization, $\bar\tau$ is represented by
$U^{-1}$. The actual frame lifts are $I$ modulo four. Of the two
possible signed lifts $\pm U^{-k}$ only the positive one can be $I$
modulo four, and it does so exactly when $k\equiv0\pmod4$.
Consequently
\[
 q_{12}(\beta_\Phi)=q_{12}(A_\Phi)=U^{-k},\qquad
 \rho_\Phi\equiv-k/2\pmod3.
\]
The second assertion follows by applying the abelianization
$\epsilon(U^2)=1$ as in FP.24. CRT now gives
$8\rho_\Phi\equiv-k\pmod{12}$. Therefore
\[
 T_\Phi=W_\Phi=I,\qquad\delta(\Phi)=0.
 \tag{FP.29.2}
\]
Taking the quotient of the genuine action homomorphism by inner
automorphisms, write
$\rho_{\mathrm{out}}:\widehat{GT}_0\to\operatorname{Out}(Q_*)$.
We have proved that $\delta$ vanishes on
$\mathcal H\cap\ker\rho_{\mathrm{out}}$. Thus it factors through
$\rho_{\mathrm{out}}(\mathcal H)$, with no choice of a finite inner
conjugator or an unproved gauge recovery.

For arbitrary multipliers, apply FP.25 with $\rho_{\mathrm{out}}$
in place of $\rho$. This gives a modulus $M$ and a finite Galois table
with the same exact evaluation formula. The multiplier residue can
also be retained in an outer finite action, as follows. The level-two
free subgroup $F\subset C_2\amalg C_3$ is characteristic: every
epimorphism onto $S_3$ takes $s$ to a transposition and $r$ to a
three-cycle, and all six such generating pairs are conjugate in $S_3$.
Thus they have the same kernel. The subgroup
\[
 F_M=\ker\bigl(F\longrightarrow F^{\mathrm{ab}}/M\bigr)
\]
is characteristic in $F$ and hence in $C_2\amalg C_3$.
In the finite quotient $(C_2\amalg C_3)/F_M$, the genuine GT action
on $F/F_M\cong(\mathbb Z/M)^2$ is the scalar $\lambda$.
Inner automorphisms act there through the six integral matrices of
the $S_3$ action on the three peripheral classes. For $M\geq2$, its
only scalar matrix is $I$; this follows directly from the transposition
$\left(\begin{smallmatrix}0&1\\1&0\end{smallmatrix}\right)$ and
the three-cycle
$\left(\begin{smallmatrix}0&-1\\1&-1\end{smallmatrix}\right)$.
Hence two genuine GT actions with equal outer residues have the same
$\lambda\bmod M$.
Adjoining this quotient, that is replacing $K_*$ by $K_*\cap F_M$,
therefore makes the full $\Delta$ a function of one finite
\emph{outer} action residue. The arithmetic modulus and Galois table
remain uncomputed; no claim of minimality is made.

The exact check of all $576$ elements of $Q$ finds
$|C_Q(\bar\tau)|=24$. The $M_2$ coinvariants exclude the twelve
noncusp centralizer residues and retain exactly the twelve cusp powers.
The three normalization cases $k=0,4,8$ give respectively
$\rho=0,1,2$ modulo three and $T=W=I$.
These checks are included in
\texttt{check\_frame\_separating\_module.py} and its output.

\paragraph{FP.30 (Proved: every induced forgetful detector annihilates the pentagon defect).}
There is an exact obstruction to obtaining the missing parity relation
by merely inducing finer four-point quotients to five points.
Write
\[
 (X_1,X_2,X_3,X_4,X_5)
   =(x_{12},x_{23},x_{34},x_{45},x_{51})
\]
for the standard cyclic boundary twists in $P_5$. In the label-five
stabilizer $B_4/\langle z_4\rangle$, with generators $a,b,c$ as in
FP.26, they have the exact words
\[
 (X_1,X_2,X_3,X_4,X_5)
       =(a^2,b^2,c^2,(aba)^2,(cbc)^2).
 \tag{FP.30.1}
\]
For $X_4,X_5$, this is the equality of the positive twist around two
points with the positive twist around the complementary three points
on a punctured sphere. It keeps the central factor until passage to
$B_4/\langle z_4\rangle$. Under the literal fold of FP.20,
\[
 q_5(X_1,X_2,X_3,X_4,X_5)=(x,y,x,1,1),
 \quad x=u^2,\quad y=v^2.
 \tag{FP.30.2}
\]
Indeed $q(a)=q(c)=u$, $q(b)=v$, and
$q((aba)^2)=q((cbc)^2)=(uvu)^2=1$ in $C_2\amalg C_3$.
The order-five automorphism $\rho$ of Lochak--Schneps (1997),
Section 1, is conjugation by $(\sigma_4\sigma_3\sigma_2\sigma_1)^2$
and sends $X_i$ to $X_{i+3}$, with indices modulo five.
Its permutation is $(1\,4\,2\,5\,3)$, so the maps
$q_5\rho^k$ run through the five forgotten labels. Their images of
the $X_i$ are exactly the five cyclic rotations of (FP.30.2), with
the specified base-path conjugations already included.

For any $f\in\widehat F_2'$ satisfying the 2-cycle relation, put
\[
 \mathcal P(f)=f(X_3,X_4)f(X_5,X_1)f(X_2,X_3)
                         f(X_4,X_5)f(X_1,X_2).
 \tag{FP.30.3}
\]
This is the actual five-cycle defect, before it is assumed trivial.
Since a commutator word is trivial when either argument is one,
\[
 q_5(\mathcal P(f))=f(y,x)f(x,y)=1.
\]
For every cyclic rotation there are again exactly two surviving
factors, one $f(x,y)$ and one $f(y,x)$; all intervening factors are
one. Hence
\[
 \mathcal P(f)\in\mathcal B_5:=
                \bigcap_{i=1}^{5}\ker(q_i|_{P_5}).
 \tag{FP.30.4}
\]
The subgroup $\mathcal B_5$ is the closed Brunnian subgroup: its
members become trivial after forgetting any marked point. The
projective quotient in $q_i$ introduces no extra pure kernel, because
$P_4$ injects into $G_4/V_4$.

Now let $\eta:L\to D$ be \emph{any} finite homomorphism factoring
through the forgetful map $q_5$, including
$\eta=\theta q_5$ from FP.23. Induction from the index-five subgroup
$L$ gives a homomorphism
\[
 G_5\longrightarrow D\wr\mathfrak S_5.
\]
For a pure element its five diagonal blocks are the conjugate
forgetful images. Equation (FP.30.4) therefore makes
$\mathcal P(f)$ the identity in this induced quotient for every
2-cycle word, even if it is nontrivial in $P_5$. Changing coset
representatives only conjugates those five identity blocks.
Thus no enlargement of the modules \emph{inside the forgotten
four-point groups}, followed by this induction alone, makes its
pentagon equation a new condition. In particular the proposed
five-block linear calculation with $E_0\wr S_5$ has an identically
zero pentagon equation. Other consistency or lifting conditions are
not asserted to be automatic.

This obstruction does \emph{not} apply automatically to FP.28:
the intersection of kernels of \emph{all} homomorphisms
$G_5\to\mathfrak S_m$ is smaller than the kernel of the one induced
representation, and need not contain $\mathcal B_5$. A genuine
five-point detector must retain an appropriate part of this Brunnian
subgroup as well as the frame data. The exact substitutions for all
five labels, the cycle permutation, and a nontrivial element of
$\mathcal B_5$ are checked in
\texttt{check\_brunnian\_detector.py}.

\paragraph{FP.31 (Proved: a small affine quotient retaining genuine Brunnian data).}
There is an explicit next quotient beyond the forgetful family.
In the literal Ferrari splitting of FP.26, reduce the kernel's
abelianization modulo three and use the signed modular action:
\[
 \eta_3:L=B_4/\langle z_4\rangle
       \longrightarrow\mathbb F_3^2\rtimes\operatorname{SL}_2(\mathbb F_3),
 \qquad
 a\longmapsto(0,U),\quad b\longmapsto(0,L_0),\quad
 c\longmapsto(e_1,U).
 \tag{FP.31.1}
\]
The target has order $9\cdot24=216$ and this map is onto.
The three Artin relations hold by direct multiplication.
Its kernel image contains both basis translations: $ca^{-1}$ maps
to $e_1$, and conjugation by $b$ supplies the second direction.
The projection of $a,b$ generates $\operatorname{SL}_2(\mathbb F_3)$,
which proves surjectivity.
The element $d=abc$ maps to $(-e_2,S)$, so
\[
 \eta_3(d)^4=\bigl(-(I+S+S^2+S^3)e_2,I\bigr)=1.
\]
Thus the map really descends through the four-strand centre and is
a representation of the five-point stabilizer, not merely of $B_4$.

For the commutator convention $[g,h]=ghg^{-1}h^{-1}$, the element
\[
 w=[[X_1,X_2],[X_3,X_4]]
\]
lies in $\mathcal B_5$: under each pattern in FP.30.2 one of the two
inner commutators is one. Nevertheless exact multiplication gives
\[
 \eta_3(w)=
 \begin{pmatrix}1&0&0\\0&1&2\\0&0&1\end{pmatrix}\ne I.
 \tag{FP.31.2}
\]
This proves both that the Brunnian subgroup is nontrivial and that
the quotient (FP.31.1) retains data absent from all five forgetful
quotients. The same affine construction modulo two sends every
$X_i$ to the identity. Hence three is the smallest coefficient
modulus in \emph{this affine family} with nontrivial pure or Brunnian
data. No assertion that order $216$ is the smallest possible finite
Brunnian quotient is made.

For a diagnostic of the equation itself, the 2-cycle word
$f_0=[x,y]$ has
\[
 \eta_3(\mathcal P(f_0))=
 \begin{pmatrix}1&0&1\\0&1&1\\0&0&1\end{pmatrix}\ne I.
 \tag{FP.31.3}
\]
Thus the pentagon is no longer an identically invisible word in this
quotient. The pair $(1,f_0)$ is \emph{not} asserted to satisfy the
3-cycle relation or define an element of $\widehat{GT}_0$; this is
a check of the comparison map, not a mixed-sign shadow or a lifting
claim.

To retain both types of data, induce the product
\[
 L\longrightarrow E_0\times
          (\mathbb F_3^2\rtimes\operatorname{SL}_2(\mathbb F_3)),
 \qquad h\longmapsto(\theta q_5(h),\eta_3(h))
\]
to $G_5$ and, if a characteristic quotient is required, apply the
all-homomorphisms construction of FP.28 to its finite coset
representation. The resulting quotient retains the actual frame
residues and the nontrivial Brunnian component (FP.31.2).
It is an explicit candidate for the next necessary pentagon equation.
The unresolved comparison is whether its equation, on actual
multiplier-one GT action residues, links the two frame parities;
neither (FP.31.2) nor (FP.31.3) proves such a link.

\paragraph{FP.32 (Proved: an integral cubic frame character and the signed-lift correction).}
Let $\mathcal H_5=\widehat{GT}\cap\mathcal H$. On this group there
are canonical continuous additive characters
$\kappa_a,\kappa_b,\kappa_d:\mathcal H_5\to\widehat{\mathbb Z}$
such that, in the ordered cyclic basis $(X_1,\ldots,X_5)$ of
$P_5^{\mathrm{ab}}\cong\widehat{\mathbb Z}^{5}$,
\[
 \begin{split}
 \alpha_a^{\mathrm{ab}}&=\kappa_a(1,1,1,1,1),\\
 \alpha_b^{\mathrm{ab}}&=\kappa_b(-1,-2,-1,1,1),\\
 \alpha_d^{\mathrm{ab}}&=\kappa_d(0,0,0,0,1),\qquad
 \kappa_b=\rho/2.
 \end{split}
 \tag{FP.32.1}
\]
Moreover the cubic CM parity has the exact five-point expression
\[
 \boxed{\ \epsilon_3(T_\Phi)=\kappa_d(\Phi)\bmod2\ },
 \qquad
 \Delta(\Phi)=\kappa_d(\Phi)+\epsilon_2(W_\Phi)\pmod2.
 \tag{FP.32.2}
\]

We give the lattice and sign proofs. The five cyclic boundary classes
form a basis because the presentation of $P_5$ on them has only
commutator relators in abelianization. The other five boundary classes
are
\[
 [x_{i,i+2}]=-[X_i]-[X_{i+1}]+[X_{i+3}],
 \quad i\pmod5.
\]
These identities follow from the five sphere relations
$\sum_{j\ne i}[x_{ij}]=0$ together with the full-twist relation;
they also verify those relations directly. Conjugation on this lattice
is the ordinary label-permutation action. For $b=(1234)$ and
$d=(234)$ its matrices are
\[
 B=\begin{pmatrix}
 0&0&0&0&-1\\1&0&1&0&0\\0&1&0&0&1\\
 0&0&-1&0&0\\0&0&-1&1&-1
 \end{pmatrix},\qquad
 D=\begin{pmatrix}
 -1&0&0&-1&0\\-1&0&-1&0&0\\0&1&-1&1&0\\
 1&0&0&0&0\\0&0&1&-1&1
 \end{pmatrix}.
 \tag{FP.32.3}
\]
The action of $a$ is the cyclic permutation of the five basis vectors.
The kernels of $B-I,D-I,a-I$ are the three rank-one lattices in
(FP.32.1). This is an integral assertion: the nonzero Smith invariants
are respectively $(1,1,1,2)$, $(1,1,1,3)$ and $(1,1,1,1)$.
The same kernels over $\widehat{\mathbb Z}$ follow because every
nonzero integer is a non-zero-divisor there.

Each conjugating word $F_i$ in FP.27 lies in $P_5'$. Hence a
multiplier-one genuine action is the identity on $G_5/P_5'$ and
$E_t^{\mathrm{ab}}=0$. Its frame must therefore lie in the indicated
fixed lattice. On $P_5^{\mathrm{ab}}$ the action is the identity,
so the crossed frame law becomes addition; uniqueness and continuity
of the frames give the three continuous characters.
The forgetful abelian map is
\[
 \begin{pmatrix}1&0&1&0&0\\0&1&0&0&0\end{pmatrix}.
\]
Its value on the $b$-fixed primitive vector is $(-2,-2)$, whereas
$\beta^{\mathrm{ab}}=(-\rho,-\rho)$ by R3. This proves
$2\kappa_b=\rho$ and the displayed normalization.

The signed projection of $\eta_3$ onto
$\operatorname{SL}_2(\mathbb F_3)$ is \emph{not} the signed lift
of the forgotten free-group frame. Write $\zeta:P_5\to\mathbb F_2$
for their difference:
\[
 Q_L(h)=(-I)^{\zeta(h)}q_3(q_5(h)),\qquad h\in P_5.
 \tag{FP.32.4}
\]
Both sides have the same projective image; their central ratio makes
$\zeta$ a continuous homomorphism. On the five generators its values
are $(0,0,0,1,1)$, because the signed images of $X_4,X_5$ are
$S^2=-I$, whereas their forgotten images are one.
Consequently $\zeta(\alpha_d)=\kappa_d\bmod2$.

By FP.26 and R3, the actual action on the affine quotient satisfies
\[
 \eta_3\Phi_4=\operatorname{Inn}(0,V_\Phi)\eta_3,
 \qquad V_\Phi=U^{8\rho_\Phi}.
\]
Here the homology matrix is exactly $V_\Phi$: it fixes $e_1$,
commutes with $U$, and has the normalized covariance with $L_0$.
There is no remaining scalar choice in this step.
The element $d=a^2bc$ has affine image
\[
 \eta_3(d)=\left((-1,-1)^t,R_0^{-1}\right),\qquad
 R_0^{-1}=\begin{pmatrix}-1&1\\-1&0\end{pmatrix}\pmod3.
\]
Its centralizer projects exactly to $\langle R_0\rangle$, the
order-three subgroup of its order-six linear centralizer. Indeed
the translation has nonzero class in
$\mathbb F_3^2/(I-R_0^{-1})\mathbb F_3^2$; the order-three subgroup
fixes that class, while multiplication by $-I$ negates it.
The affine frame normalized by $(0,V_\Phi)$ centralizes this element.
Its linear projection is
$V_\Phi^{-1}(-I)^{\zeta(\alpha_d)}q_3(A_\Phi)$.
It lies in the order-three subgroup, so applying the unique sign
character of the order-six centralizer gives
$\epsilon_3(T_\Phi)=\zeta(\alpha_d)$, which proves (FP.32.2).
The exact centralizer has order nine and its linear projection has
order three; both are independently enumerated in the attached check.

The integral lattice alone supplies no equality with the harmonic
parity. In $G_5/P_5'$ the identity automorphism admits independently
every one of the three invariant frame vectors in (FP.32.1). Their
conjugations fix $a,b,d$ in that quotient, so all presentation relations
hold for every choice of coefficients. In particular the abelianization
of (FP.27.5) reduces only to $D\alpha_d^{\mathrm{ab}}
=\alpha_d^{\mathrm{ab}}$; it imposes no relation between
$\kappa_b$ and $\kappa_d$. These are abelianized primitive choices,
not actual GT frames or a lifting assertion. The remaining genuine
comparison is precisely
\[
 \epsilon_2(W_\Phi)\stackrel{?}{=}\kappa_d(\Phi)\pmod2
 \quad(\Phi\in\mathcal H_5).
\]
On Galois this follows from R4 and equals the specified common
quadratic character. Identifying the entire integral $\kappa_d$ with
a normalized half-Kummer character would require its own path and
normalization proof; only the parity identity is asserted here.
The lattice, signed correction, and both affine centralizers are
checked in \texttt{check\_five\_point\_lattice.py} and its output.

\paragraph{FP.33 (Proved: the exact affine pentagon equation and its limit).}
For $\Phi=(1,f)\in\mathcal H$, let
\[
 D_6=\langle p,t\mid p^3=t^2=(pt)^2=1\rangle\cong S_3,
 \qquad f(p,t)=p^{a_f},\quad a_f\in\mathbb F_3.
 \tag{FP.33.1}
\]
The exponent exists uniquely because $f\in\widehat F_2'$ and
$D_6'=\langle p\rangle$. Then the exact value of the cyclic
pentagon defect in the affine quotient of FP.31 is
\[
 \boxed{\ \eta_3(\mathcal P(f))
             =\bigl((a_f-\rho_\Phi)e_1,I\bigr)\ }.
 \tag{FP.33.2}
\]
Consequently every genuine multiplier-one GT element satisfies the
necessary equation
\[
 a_f=\rho_\Phi\pmod3.
 \tag{FP.33.3}
\]
The difference $\eta_{\mathrm{aff}}(\Phi)=a_f-\rho_\Phi$ is a
continuous additive character $\mathcal H\to\mathbb F_3$.
This is an actual equation for globally defined pairs; finite
admissibility or an unspecified lift is not used.

Here is a complete calculation of (FP.33.2).
In $V\rtimes\operatorname{SL}_2(\mathbb F_3)$, $V=\mathbb F_3^2$,
the five pure generators have images
\[
 (0,U^2),\quad(0,L_0^2),\quad(-e_1,U^2),\quad
 (0,-I),\quad(e_2,-I).
 \tag{FP.33.4}
\]
Set $p=(e_1,I)$, $t=(0,-I)$ and write the same letters $X_i$
for these five images. The element $h_0=(0,U^2)$ commutes with
$p,t$, and $X_3=p^{-1}h_0$. Inserting a central argument in a
commutator word changes nothing. Conjugation by $t$ inverts $p$,
so
\[
 f(X_3,X_4)=f(p^{-1},t)=p^{-a_f}.
\]
For $t'=X_5$, one has $t'pt'^{-1}=p^{-1}$ and
$t'X_1t'^{-1}=pX_1$. Thus $h=p^2X_1$ commutes with $p,t'$,
and $X_1=ph$. A translation in the $e_2$ direction conjugates
$t$ to $t'$ while fixing $p$ and every element of $\langle p\rangle$.
The 2-cycle relation gives
\[
 f(X_5,X_1)=f(t',p)=f(p,t')^{-1}=p^{-a_f}.
\]
The pair $(X_2,X_3)$ is the conjugate, by the translation $-e_2$,
of the linear pair $(L_0^2,U^2)$. If
$F=f(U^2,L_0^2)$, its contribution is therefore
\[
 f(X_2,X_3)=\bigl((F^{-1}-I)e_2,F^{-1}\bigr).
\]
For the fourth factor put $q=(e_2,I)$, so $X_4=t$ and $X_5=qt$.
The multiplier-one 3-cycle relation at $(x,y,z)=(t,qt,q)$ is
\[
 f(q,t)f(qt,q)f(t,qt)=1.
\]
Conjugating $(t,q)$ by $q^2$ gives $(qt,q)$; the value of a
commutator word is in $\langle q\rangle$ and is unaffected by this
conjugation. Hence $f(qt,q)=f(t,q)=f(q,t)^{-1}$, and
$f(X_4,X_5)=1$. Finally $f(X_1,X_2)=(0,F)$.
Multiplying in the order of (FP.30.3) gives
\[
 \eta_3(\mathcal P(f))
       =\bigl(-2a_f e_1+(F^{-1}-I)e_2,I\bigr).
\]
R3 supplies
\[
 F=\begin{pmatrix}1&\rho\\\rho&1+\rho^2\end{pmatrix}\pmod3,
 \qquad (F^{-1}-I)e_2=-\rho e_1,
\]
which is exactly (FP.33.2).

To verify the character assertion, the standard pure automorphism of
the pair descends under $x\mapsto p,y\mapsto t$ to
\[
 p\longmapsto p,\qquad
 t\longmapsto p^{-a_f}tp^{a_f}.
\]
It is an inner automorphism of $D_6$ and is the identity on its
derived group. The GT composition formula for $f$ therefore adds
the two $a_f$ exponents. The harmonic coordinate $\rho$ is additive
on $\mathcal H$ by the abelianized frame law, proving the claim.

At affine modulus four, every matrix of a generator $X_i$ is
congruent to $I$ modulo two. The group
$1+2\operatorname{Mat}_3(\mathbb Z/4)$ is abelian. Hence every
$f(X_i,X_j)$ is already one for $f\in\widehat F_2'$, and the
pentagon has no further component there. CRT shows that the entire
affine modulus-twelve equation is precisely (FP.33.3).

This necessary equation cannot force the CM parity equality, even on
actual globally liftable elements of $\widehat{GT}_0$.
Let $g\in\mathcal H$ be the verified R2 element, so
$\Delta(g)=1$. Its \emph{GT-group cube} $g^3$ satisfies
\[
 \eta_{\mathrm{aff}}(g^3)=3\eta_{\mathrm{aff}}(g)=0,
 \qquad \Delta(g^3)=3\Delta(g)=1.
 \tag{FP.33.5}
\]
Thus its pair satisfies the full affine modulus-twelve pentagon
equation and has mixed CM parity. This is the cube in the group of
automorphisms, not the literal substitution $f\mapsto f^3$.
It is not claimed to satisfy the actual five-point pentagon, nor to
define an automorphism of every finite five-point quotient. What it
proves is the exact failure of this necessary finite equation to
imply $\Delta=0$, with an already global $\widehat{GT}_0$ witness.
The next comparison must retain additional Brunnian information
which is not reduced to this order-three character.
All finite affine identities used in the derivation are checked in
\texttt{check\_brunnian\_detector.py}; no pentagon-shadow census is
used as a substitute for a lifting theorem.

\paragraph{FP.34 (Proved: the harmonic parity is an integral diagonal beta phase).}
There is an exact finite mixed-coefficient formula for the other
parity. Set
\[
 R=(\mathbb Z/4)[\omega]/(\omega^2+\omega+1),\qquad
 c=(1-\omega)^{-1}=2+3\omega.
\]
In the finite affine group $R_{\rm add}\rtimes\langle-\omega\rangle$
define
\[
 s\longmapsto(c,-1),\qquad r\longmapsto(c,\omega).
 \tag{FP.34.1}
\]
Their orders are two and three. Thus this is a continuous
homomorphism from $\widehat B=C_2\amalg C_3$; its image has order
$96$. The exact peripheral images are
\[
 \tau=sr\longmapsto(0,-\omega),\qquad
 x=(sr)^{-2}\longmapsto(0,\omega),\qquad
 y=sxs\longmapsto(1,\omega).
\]
For $\Phi=(1,f)\in\mathcal H$, write the image of $f$ as
$(k_f,1)$ and set
\[
 \mathcal B_f=1-(1-\omega)k_f\in R.
\]
Then
\[
 \boxed{\ \mathcal B_f
       =\omega^{\rho_\Phi}(-1)^{\epsilon_2(W_\Phi)}\quad\text{in }R\ }.
 \tag{FP.34.2}
\]
This formula retains the coefficient-two component together with
the order-three eigenvalue, and does not pass to a pro-two quotient.

To prove it, write the image of the actual harmonic frame as
$\beta=(b,\omega^\rho)$. Its linear component is as stated because
$\beta^{\mathrm{ab}}=(-\rho,-\rho)$, so it is
$\omega^{-2\rho}=\omega^\rho$. The exact equality
$\beta s\beta^{-1}=sf$ becomes
\[
 2b+\omega^\rho c=c-k_f.
 \tag{FP.34.3}
\]
The image of the shear $V=U^{8\rho}$ is $(0,\omega^\rho)$,
because $U=\tau^{-1}$ has zero translation and linear component
$-\omega^{-1}$. Hence
\[
 W=\beta^{-1}V\longmapsto(w,1),\qquad
 w=-\omega^{-\rho}b.
\]
Modulo two, (FP.34.1) is the projective level-three quotient
$A_4=\mathbb F_4\rtimes C_3$ with the prescribed images of
$x,y,s,r$. The corrected harmonic frame belongs to the cyclic
centralizer of the signed $S$; its projective image is therefore
either $1$ or $s$. Accordingly
\[
 w/c\bmod2=\epsilon_2(W)\in\{0,1\}\subset\mathbb F_4.
\]
Multiplying (FP.34.3) by $c^{-1}$ yields
$2b/c=\mathcal B_f-\omega^\rho$, and hence
\[
 \mathcal B_f=\omega^\rho(1-2w/c)
             =\omega^\rho(-1)^{\epsilon_2(W)}\pmod4.
\]
This proves (FP.34.2), including the division by two which would
be lost by an initial reduction modulo two.

In the usual metabelian notation
$f\equiv A_f(x,y)[x,y]\pmod{F''}$, with
$[x,y]=xyx^{-1}y^{-1}$, the commutator image in (FP.34.1) is
$\omega-1$. Therefore
\[
 k_f=(\omega-1)A_f(\omega,\omega),\qquad
 \mathcal B_f=1+(\omega-1)^2A_f(\omega,\omega)
            =1-3\omega A_f(\omega,\omega).
 \tag{FP.34.4}
\]
The plus sign in this intrinsic multiplier-one beta normalization
agrees with the opened primary proof of
Nakamura--Tsunogai--Yasuda (2010), Proposition 4.4, printed p.~443.
Identifying its exact labelled value with Ihara's
$B_\Phi(1/3,1/3)$ additionally requires matching the free-group
marking and the action convention. Simultaneous inversion of the
roots of unity, or inversion of the resulting beta value, leaves
its sign unchanged but can change the displayed power of $\omega$.
No distribution formula with its cyclotomic exponent is imported
without that comparison. The intrinsic formula (FP.34.2) is
independent of this naming issue.

For genuine GT elements, combining (FP.32.2) and (FP.34.2) reduces
the main sign equality to the precise finite identity
\[
 \mathcal B_f\stackrel{?}{=}
            \omega^{\rho_\Phi}(-1)^{\kappa_d(\Phi)}
 \quad\text{in }(\mathbb Z/4)[\omega].
 \tag{FP.34.5}
\]
Full profinite beta/gamma factorization supplies
$B(s,t)=\Gamma(s)\Gamma(t)/\Gamma(s+t)$, but that formula alone
does not identify the finite phase in (FP.34.5) with a Kummer or
cubic-frame coordinate. Ihara's further multiplication conditions
must be distinguished from the proved factorization; their use
would require a separate theorem valid on all $\widehat{GT}$.
The exact ring calculations, group order, peripheral identities,
and all six $(\rho\bmod3,\epsilon_2)$ phase cases are checked in
\texttt{check\_harmonic\_beta\_phase.py} and its output.

\paragraph{FP.35 (Proved: a finite solvable detector of derived length at most four).}
The normalized finite detector of (FP.24), and the outer detector
of (FP.29), can be chosen as actions on a finite characteristic
quotient of $\widehat B$ whose derived length is at most four.
On $\mathcal H$ this detects $\delta$ directly. For all of
$\widehat{GT}_0$, use exactly the finite multiplier residue and
Galois normalization table in (FP.25); neither an explicit value
of its modulus $M$ nor a smaller source quotient is asserted here.

Indeed, the finite group
$Q=\operatorname{SL}_2(\mathbb Z/12)/\{\pm I\}$ has derived
series of orders
\[
                     576,\quad96,\quad8,\quad1.
\]
For comparison the local series have orders
$48,12,4,1$ at four and $24,8,2,1$ at three. CRT gives the
first displayed series: the simultaneous central element
$(-I,-I)$ is not in the first derived product, because
$-I\notin\operatorname{SL}_2(\mathbb Z/4)'$.
These are independently enumerated in
\texttt{check\_solvable\_detector.py} and its output.

The group $E=(M_2\oplus M_3)\rtimes Q$ of (FP.23) is an
extension of a group of derived length three by an abelian group.
Its fourth derived group is therefore trivial. The characteristic
quotient $Q^*$ of (FP.24) embeds in the finite product of the
images of all homomorphisms $\widehat B\to E$. Subgroups and
finite products preserve the bound on derived length, so
$(Q^*)^{(4)}=1$.

For the full multiplier version of the outer detector, take the
joint characteristic quotient with $\widehat B/F_M$ as in
(FP.29). That added factor is an extension of $S_3$ by the
abelian group $F/F_M\simeq(\mathbb Z/M)^2$, hence has derived
length at most three. The joint source quotient still has derived
length at most four. The conclusion concerns the finite
\emph{source group} on which the automorphisms act; it does not
assert the same bound for its automorphism group, nor assert that
every finite automorphism or pentagon assignment lifts globally.

\paragraph{FP.36 (Proved, using the actual strengthened R2 lift): a simultaneous obstruction to low-derived and two-adic Burau pentagon certificates.}
The lifting audit's strengthened LI.8 supplies an \emph{actual}
$k=(1,f_k)\in\mathcal H$ satisfying
\[
 \Delta(k)=1,\qquad
 f_k,\beta_k\in R_{\rm nil}(F'')\subset F''',\qquad A_k\in F''.
 \tag{FP.36.1}
\]
Here the derived groups and the nilpotent residual are closed.
The strengthened residual assertion follows from its compactness
construction as follows. In the profinite quotient
$\widehat B/R_{\rm nil}(F'')$, let $q$ be any finite quotient
and put $M=F''\ker(q)$. This is an open normal subgroup of
$\widehat B$, and $q(M)=q(F'')$ is nilpotent. Consequently
$q$ kills $R_{\rm nil}(M)$. Taking all such $q$ proves
\[
 \bigcap_{M\mathrel{\triangleleft_o}\widehat B,\ F''\subset M}
             R_{\rm nil}(M)=R_{\rm nil}(F'').
\]
The reverse inclusion follows by restriction to $F''$ of each
finite nilpotent quotient of $M$. Finally,
$R_{\rm nil}(F'')\subset(F'')'=F'''$ because finite abelian
quotients are nilpotent.

\emph{Simultaneous pure-derived-length-three consequence.}
In every finite quotient of $G_5$ whose pure image has derived
length at most three, all five evaluations
$f_k(X_i,X_j)$ in the pentagon are one. Hence its pentagon
word is one in every such quotient, simultaneously, although
$\Delta(k)=1$. More generally, every standard Drinfeld
generator formula at every number of marked points has the form
of a half twist conjugated by an evaluation of $f_k$ at two pure
elements (with multiplier one). It evaluates to that same half
twist modulo $P_n'''$. Thus these formulas give the literal
identity automorphism on all $G_n/P_n'''$ simultaneously.
This does not assert an action on $G_n$ itself, or membership of
$k$ in $\widehat{GT}$. It proves that no certificate confined to
these quotients can force the CM sign equality. The witness is
the already lifted profinite element (FP.36.1), not a collection
of unrelated finite solutions.

\emph{A genuinely five-point quotient, with its central relation checked.}
Let $R_n=(\mathbb Z/2^n)[\omega]/(\omega^2+\omega+1)$ and set
\[
 A=\begin{pmatrix}1-\omega&\omega&0\\1&0&0\\0&0&1\end{pmatrix},
 \quad B=\begin{pmatrix}1&0&0\\0&1-\omega&\omega\\0&1&0\end{pmatrix},
 \quad C=\begin{pmatrix}1&0&-1\\0&1&-1\\0&0&-\omega\end{pmatrix}.
 \tag{FP.36.2}
\]
These reduced Burau matrices satisfy the three braid and distant
commutation relations. Before projectivization they satisfy
$(ABC)^4=\omega I$. Therefore their projective classes give an
actual continuous quotient of the label-five stabilizer
$L=B_4/\langle z_4\rangle$.
Inducing this homomorphism from the actual index-five subgroup
$L\subset G_5$ gives a finite quotient of $G_5$ in a wreath
product. In particular no unverified direct representation of
$B_5$ through the spherical relator is being used.

For $n=2$, the projective image has order $55\,296$; its
modulo-two image has order $216$ and its reduction kernel is
$(\mathbb F_2)^8$. This is genuinely finer than the earlier
affine quotient. Nevertheless it still cannot force $\Delta=0$.
In pentagon order, the five evaluation pair images have derived
series
\[
\begin{array}{c|c|c}
 \text{pair}&\text{modulus }4&\text{modulus }8\\ \hline
 (X_3,X_4),(X_5,X_1)&144,12,4,1&1152,96,32,2,1\\
 (X_2,X_3),(X_1,X_2)&192,64,4,1&1536,512,8,1\\
 (X_4,X_5)&48,12,4,1&384,96,32,2,1 .
\end{array}
 \tag{FP.36.3}
\]
Thus the modulus-four equation is killed already by
$f_k\in F'''$. The modulus-eight groups do reach derived length
four, but (FP.36.1) still kills every evaluation: the second
derived group of each pair image is a finite two-group, hence
nilpotent.

In fact this obstruction holds at \emph{every} $2^n$ in the
same family. At modulus two the pair images have series
$18,3,1$, $24,8,2,1$, and $6,3,1$, respectively.
Their second derived groups are trivial or $C_2$. The kernel
of reduction from $\operatorname{PGL}_3(R_n)$ to
$\operatorname{PGL}_3(\mathbb F_4)$ is a two-group: its
successive congruence quotients are additive matrix groups in
characteristic two, modulo their scalar subspaces. Hence the
second derived group of every lifted pair image is a two-group.
The image of $F''$ under that evaluation lies in this nilpotent
group, so $R_{\rm nil}(F'')$ maps to one. Therefore
\[
           \eta_{2^n}\bigl(\mathcal P(f_k)\bigr)=1
           \quad\text{for every }n\geq1,
           \qquad \Delta(k)=1.
 \tag{FP.36.4}
\]
The same conclusion holds after induction to $G_5$ in every
block. Choose the five coset representatives to be the powers
of the transitive rotation used in (FP.30). Its conjugation
cyclically permutes the five $X_i$, so the pairs in each block
are again the five cyclically adjacent pairs listed in
(FP.36.3); each corresponding $f_k$ value is one.
The program \texttt{check\_burau\_pair\_filtration.py} and
the three outputs \texttt{check\_burau\_pair\_mod2.out},
\texttt{check\_burau\_pair\_mod4.out}, and
\texttt{check\_burau\_pair\_mod8.out} verify the literal
matrices, central word, and all finite series used here.
This failure leaves open quotients whose evaluation pair has a
non-nilpotent second derived image; it makes no assertion that
the full pentagon is satisfied by $k$.

\paragraph{FP.37 (Proved: an exact quadratic necessary pentagon character).}
There is a continuous homomorphism
$\nu_4:\mathcal H\to\mathbb F_2$, defined by one literal
$S_4$ evaluation, such that
\[
 \boxed{\quad \epsilon_2(W_\Phi)=\nu_4(\Phi)
                 \qquad(\Phi\in\widehat{GT}\cap\mathcal H).\quad}
 \tag{FP.37.1}
\]
Thus $\epsilon_2(W)+\nu_4$ is a finite continuous character
of the actual gentle group which vanishes on its genuine GT
subgroup. Its nonvanishing on the gentle group is not asserted.
It is not identified here with $\Delta$, with the cubic-frame
character, or with the half-Kummer character.

\emph{Definition and homomorphism law.}
In $S_4$ put $p=(123)$ and $t=(1234)$, with products acting
right to left. The orders of $p,t,pt$ are $3,4,4$.
For $\Phi=(1,f)\in\mathcal H$, write $R=f(p,t)\in A_4$.
The actual peripheral automorphism of the free group preserves
the conjugacy class of $xy$, hence $pR^{-1}tR$ has order four.
Writing $A_4=V_4\rtimes\langle p\rangle$, the four choices
of $v\in V_4$ give order four for $pv^{-1}tv$ exactly when
$v\in\{1,t^2\}$; the other two choices give order two.
Consequently there are unique $\nu_4\in\mathbb F_2$ and
$a\in\mathbb F_3$ with
\[
                 f(p,t)=t^{2\nu_4}p^a.
 \tag{FP.37.2}
\]
The quotient $S_4/V_4=S_3$ identifies $a$ with the actual
$a_f$ of (FP.33). Moreover the induced action on this $S_4$
quotient is exactly $\operatorname{Inn}(p^{-a_f})$:
it fixes $p$ and sends $t$ to $p^{-a_f}tp^{a_f}$.
In particular its kernel is preserved by every actual gentle
element; no invariance of a Burau kernel has been assumed.
For two such elements the GT multiplication rule gives
\[
 R_{\Phi\Psi}
   =R_\Phi p^{-a_\Phi}R_\Psi p^{a_\Phi}
   =t^{2(\nu_4(\Phi)+\nu_4(\Psi))}
                  p^{a_\Phi+a_\Psi}.
\]
This proves the homomorphism assertion and continuity follows
from the finite evaluation.

\emph{Reduction of the actual pentagon.}
Use the modulus-four matrices in (FP.36.2), with
$X_1=A^2,X_2=B^2,X_3=C^2,X_4=(ABA)^2,X_5=(CBC)^2$.
Let $D$ be the image of the reversing half twist
$abcaba$. Then $D^2=1$, and its conjugation interchanges
$X_1,X_3$ and $X_4,X_5$, while fixing $X_2$.
Put $F=f(X_1,X_2)$ and $R=f(X_3,X_4)$. The two-cycle
therefore gives the exact matrix pentagon
\[
 \mathcal P(f)=
 R(DR^{-1}D)(DF^{-1}D)\,f(X_4,X_5)\,F.
 \tag{FP.37.3}
\]
Inside $\langle X_3,X_4\rangle$ one can strip a central
element $h$ of order three from $X_3$, obtaining
$p_0=X_3h^{-1}$ and $t_0=X_4$ which generate a group of
order $48$. Its centre is generated by an involution $e_0$;
its quotient by that centre is the marked $S_4$ of
(FP.37.2). Its derived subgroup is an $A_4$ mapping
isomorphically onto the quotient's derived subgroup, so the
commutator evaluation is exactly
\[
 R=(e_0t_0^2)^{\nu_4}p_0^{a_f}.
 \tag{FP.37.4}
\]
Here $h$ disappears from every evaluation of $f\in F'$.

For completeness all conjugacies needed for the three-cycle
reduction are literal finite matrices. Set
$z=(X_4X_5)^{-1}$. There is an invertible matrix $C_0$ with
$C_0p_0C_0^{-1}=z$ and $C_0t_0C_0^{-1}=X_4$.
Furthermore $z^{-1}X_4z$ differs from $X_5$ by the central
involution of $\langle X_4,X_5\rangle$. Put
$H=C_0RC_0^{-1}=f(z,X_4)$. The multiplier-one three-cycle,
and the fact that central factors disappear from a commutator
word, give
\[
 f(X_5,z)=z^{-1}H^{-1}z,
 \qquad f(X_4,X_5)=z^{-1}HzH^{-1}.
 \tag{FP.37.5}
\]
Equations (FP.37.3)--(FP.37.5) reduce the test to the actual
two parameters $F$ and $R$ of one and the same $f$.

Here is an exact small certificate for the remaining matrix
identity. Encode $a+b\omega\in R_2$ by the integer $a+4b$
with $0\leq a,b<4$, and list matrices in row-major order.
The constants can be chosen as
\[
\begin{split}
 h&=(8,1,0;15,14,0;0,0,9),\\
 e_0&=(1,10,0;8,3,0;0,0,11),\\
 C_0&=(1,4,9;10,4,0;1,13,0).
\end{split}
\]
The group $H_{12}=\langle X_1,X_2\rangle$ has order $192$.
Its positive level-three map sends $X_1$ to $U^2$ and $X_2$
to $L_0^2$. Enumerate the images $\beta\in H_{12}$ of the
\emph{actual} harmonic frame, with
$q_3(\beta)^{-1}U^{8\rho}\in\langle S\rangle$.
The equation $F=(ABA)^{-1}\beta(ABA)\beta^{-1}$ gives
exactly four possible matrices $F$ for each $\rho\bmod3$,
two for each $\epsilon=\epsilon_2(W)$. They are
\[
                   F=F^0_{\rho,\epsilon}Z^c,
             \qquad Z=(ABA)^4,\quad c\in\mathbb F_2,
\]
where the six base matrices are
\[
\begin{array}{c|c|l}
 \rho&\epsilon&F^0_{\rho,\epsilon}\\ \hline
0&0&(1,0,0;0,1,0;0,0,1)\\
0&1&(1,10,0;8,9,10;0,8,3)\\
1&0&(1,7,0;10,0,14;3,10,15)\\
1&1&(1,13,0;0,0,14;9,0,5)\\
2&0&(8,1,0;4,9,12;7,12,10)\\
2&1&(0,1,0;14,3,4;15,6,0).
\end{array}
\]
For a genuine pentagon, the independently proved necessary
equation (FP.33) gives $a_f=\rho$.
Substitution in (FP.37.3)--(FP.37.5)
now yields, for all the remaining $24$ cases, the exact formula
\[
                  \mathcal P(f)=v^{\epsilon+\nu_4}w_\rho^c,
 \tag{FP.37.6}
\]
where $v,w_\rho$ are independent commuting involutions, with
encoded matrices
\[
\begin{split}
 v&=(1,10,8;8,3,8;0,0,9),\\
 w_0&=(1,8,10;10,9,8;10,2,3),\\
 w_1&=(1,2,8;2,9,2;10,2,3),\\
 w_2&=(1,8,2;2,9,10;10,2,1).
\end{split}
\]
Thus $\mathcal P(f)=1$ forces $\epsilon=\nu_4$, proving
(FP.37.1). The program
\path{check_burau_quadratic_character.py} and its output
check every conjugacy, the marked $S_4$ quotient, the 36
coordinate multiplication laws, and all 24 instances of
(FP.37.6). This is a necessary equation for actual GT elements;
none of the enumerated parameter combinations is asserted to
lift to an actual gentle element merely because it passes the
finite equations.

\emph{Exact limitation of the cubic-frame comparison in this quotient.}
The same Burau target cannot identify $\nu_4$ with
$\kappa_d\bmod2$ by the cubic conjugator equation alone.
Its entire pure image $G$ is generated by $X_1,X_2,X_3$, all
of order three: $X_4\in\langle X_1,X_2\rangle$ and
$X_5\in\langle X_2,X_3\rangle$. Hence $G$ has no
nontrivial elementary-two quotient. Goursat's lemma then makes
the joint pure map
\[
            P_5\longrightarrow G\times(P_5^{\rm ab}/2)
                              =G\times\mathbb F_2^5
\]
surjective. Choose an actual source element $h\in P_5$ with
Burau image one and abelianization $e_5$ modulo two.
There is a literal homomorphism from the order-$216$ reduction
of the Burau image to $\operatorname{SL}_2(\mathbb F_3)$
sending $A,C$ to $U$ and $B$ to $L_0$. The certificate checks
all generator multiplications on all $216$ elements and
surjectivity onto its $24$-element target. This is precisely
the signed $Q_L$ representation in (FP.32), since the three
generators have the prescribed images. Thus Burau$(h)=1$
implies $Q_L(h)=I$. The signed identity in (FP.32) gives
$q_3(q_5(h))=-I$, because $\zeta(h)=1$;
the abelian forgetting map gives $q_4(q_5(h))=I$.
Thus $q_{12}(q_5(h))=[5I]$.
Its image centralizes $d$ in the joint quotient retaining
$G$, $P_5^{\rm ab}/2$, and $q_{12}q_5$: its three coordinates
are respectively $1$, the $d$-fixed vector $e_5$, and a central
scalar. Multiplication by this finite centralizer element
changes the candidate cubic-frame parity by one while preserving
all those finite conjugator equations. This proves a precise
failure of parity extraction from this target. It does not
produce a nontrivial element of the actual profinite centralizer
$C_{P_5}(d)$, which is trivial, and it does not refute an eventual
comparison using additional five-point quotients or a new theorem.

\paragraph{FP.38 (Proved: the first relative characteristic-three homology extension is still insufficient).}
There is a concrete further test which adds characteristic-three
homology to the genuinely five-point quotient of order $216$;
it also fails to force $\Delta=0$, for an exact structural reason.
Let $L=B_4/\langle z_4\rangle$ and let
\[
 q:L\twoheadrightarrow Q=\mathbb F_3^2\rtimes
                 \operatorname{SL}_2(\mathbb F_3),\qquad |Q|=216,
\]
send $a,b,c$ respectively to $(0,U),(0,L_0),(e_1,U)$.
Put $N=\ker(q)$. Use the presentation with three generators
and relators
\[
 aba(bab)^{-1},\quad bcb(cbc)^{-1},\quad aca^{-1}c^{-1},
                  \quad(abc)^4.
\]
The Fox derivatives of these four relators and all their
$Q$-translates form an $864\times648$ matrix over $\mathbb F_3$.
The exact rank is $422$. Therefore the quotient of
$\mathbb F_3[Q]^3$ by its row module is a $226$-dimensional
$Q$-module $W$. The literal assignments
\[
       a\longmapsto(e_a,q(a)),\quad
       b\longmapsto(e_b,q(b)),\quad
       c\longmapsto(e_c,q(c))
           \qquad\text{in } W\rtimes Q
 \tag{FP.38.1}
\]
kill every defining relator and give an actual finite quotient
of $L$. They realize the usual relative abelian extension
$L/[N,N]N^3$ as their image. The boundary map to
$\mathbb F_3[Q]$ has rank $215$, so
\[
                    \dim_{\mathbb F_3}H_1(N,\mathbb F_3)=11.
\]
Neither this construction nor the pentagon evaluation assumes
that an arbitrary gentle action preserves $N$.

Here is the decisive calculation on the evaluation subgroups.
For a pair image $P$, write $H=q(P)$ and
$U_P=P\cap W$. Schreier rewriting on the finite Cayley graph
of $H$ computes $U_P$ exactly as the span of
$v_h+h v_g-v_{hg}$, for the two generators $g$ and the chosen
lifts $(v_h,h)$. It gives
\[
\begin{array}{c|c|c|c|c}
 \text{pair}&(|H|,|H'|,|H''|)&\dim U_P&|P|
       &H''\text{ action on }U_P\\ \hline
 (X_1,X_2),(X_2,X_3)&(24,8,2)&4&1944&\text{trivial}\\
 (X_3,X_4),(X_5,X_1)&(18,3,1)&5&4374&\text{trivial}\\
 (X_4,X_5)&(6,3,1)&4&486&\text{trivial}.
\end{array}
 \tag{FP.38.2}
\]
The reversing half twist gives the paired equalities in the
first two rows. The action in the last column is checked on
an exact basis of the entire kernel $U_P$, not just on a
selected commutator subspace.

Each $U_P$ is elementary abelian of characteristic three.
If $H''=1$, then $P''\subset U_P$ is abelian. In the remaining
case, the preimage of $H''=C_2$ is a central extension of a
cyclic group by $U_P$, because the displayed action is trivial;
that preimage is abelian as well. Thus every pair $P$ has
derived length at most three. The actual element $k$ in
(FP.36.1) consequently has $f_k(X_i,X_j)=1$ in all five pair
images. Its whole pentagon is one in this quotient and in the
induced five-point quotient, while $\Delta(k)=1$.

The exact rank, Fox-boundary identities, all affine relators,
Schreier spans and triviality of the entire $C_2$ action are
checked by \texttt{check\_fox\_relative\_module.py} and its
output. The complete quotient map, the $216$ module actions,
and the finite base multiplication are stored in
\texttt{fox\_relative\_module.npz}. This resolves this specific
relative homology test negatively; it does not rule out a
different finite base or a further nonabelian relative extension.

\paragraph{FP.39 (Proved, with R2): one actual mixed-parity element passes the entire Eisenstein--Burau congruence family.}
The odd-characteristic test does escape the pairwise nilpotence
criterion, but it too admits a stronger, actual lifting obstruction.
At $\omega=2\in\mathbb F_7$, the precise pair series are
\[
\begin{array}{c|c}
 (X_3,X_4),(X_5,X_1)&1008,336,336\\
 (X_1,X_2),(X_2,X_3)&1029,343,7,1\\
 (X_4,X_5)&336,336.
\end{array}
\]
The repeated terminal groups of order $336$ are perfect. Their
markings are checked explicitly: $\langle X_4,X_5\rangle$
is $\operatorname{SL}_2(\mathbb F_7)$ with
$X_4\mapsto U$, $X_5\mapsto L_0^3$;
$\langle X_3,X_4\rangle$ is
$C_3\times\operatorname{SL}_2(\mathbb F_7)$ with
$X_3\mapsto(1,(UL_0^3)^{-1})$, $X_4\mapsto(0,U)$.
All $336$ and $1008$ generator multiplications are checked in
\texttt{check\_burau\_odd\_pair.py} and its output.
No perfectness claim is made for the harmonic pair
$\langle X_1,X_2\rangle$.

The stronger whole-family conclusion follows from integral
identities over $\mathbb Z[\omega]$, not from extrapolation of
these finite computations. For the matrices of (FP.36.2),
\[
\begin{split}
 X_4-I&=(1-\omega)(1,1,1)^{\mathsf T}(1,\omega,\omega^2),\\
 X_5-I&=\begin{pmatrix}0&\omega-1&-2\omega-1\\0&0&0\\0&0&0\end{pmatrix}.
\end{split}
\]
Both square to zero. Consequently $X_4$ and $X_5$ have orders
dividing the coefficient modulus. Also $A,B$ fix the vector
$v=(1,1,1)^{\mathsf T}$ and covector
$\ell=(1,\omega,\omega^2)$, with $\ell(v)=0$.
The basis $(v,(-\omega,1,0)^{\mathsf T},e_1)$ has determinant
$-1$, so this is an integral primitive flag at every modulus.
The harmonic pair is triangular in that basis. Its first
derived image is upper unitriangular, its second derived image
is abelian, and its third derived image is one, at every finite
coefficient ring. The reversed harmonic pair has the same property.

For a finite collection of coefficient moduli away from three,
choose $m$ to be their least common multiple with two. Then
$m$ is even and prime to three. In the source free group put
\[
             A_m=[[x^m,y^m],[y^m,z^m]],
                      \qquad z=(xy)^{-1}.
 \tag{FP.39.1}
\]
This belongs to $F''$. Its literal signed level-twelve image is
$5I$ for every such $m$: the four even, prime-to-three residue
classes $2,4,8,10$ modulo twelve are all checked exactly.
For each of the three nonharmonic pair evaluations
$(X_3,X_4),(X_5,X_1),(X_4,X_5)$, at least one peripheral
is $X_4$ or $X_5$. Its $m$th power is one, so one of the two
inner commutators in (FP.39.1) is one and the outer commutator
is one. The same argument applies to all six peripheral
permutations: that peripheral is merely moved or conjugated.
Thus $A_m$ lies in the $\widehat B$-normal core of all the
specified finite evaluation kernels. Since $m$ is even, it
also lies in the section-47 cores used by the lifting audit.

R2 can now be applied exactly as in the strengthened LI.8,
intersecting its finite normal bases with these evaluation
kernels; the specified seed $A_m$ retains the scalar $5I$.
Compactness over the finite collections and over all open normal
subgroups containing $F''$ gives one actual gentle element
$k=(1,f_k)$ with
\[
 \Delta(k)=1,\qquad D(k)=1,
 \qquad f_k,\beta_k\in R_{\rm nil}(F''),
 \qquad q_{12}(A_k)=5I,
\]
whose nonharmonic pentagon evaluations are one at every modulus
away from three. The harmonic evaluations are already killed by
$f_k\in F'''$. At the ramified prime above three, all five
$X_i$ reduce to one modulo $(1-\omega)$, so their finite
congruence images are three-groups; $R_{\rm nil}(F'')$ kills
these evaluations as well. CRT proves that all five factors
are one at every finite quotient
$(\mathbb Z/N)[\omega]/(\omega^2+\omega+1)$, simultaneously.
The rotated blocks in the induced five-point quotients are
covered exactly as in (FP.36).

The complete R2/compactness argument is also recorded as LI.9
by the lifting audit. Independent integral verification is in
\texttt{check\_burau\_integral\_identities.py} and its output.
This is a theorem about an actual global gentle witness passing
an entire congruence family. It neither places that witness in
GT nor identifies the full profinite five-point relation with
any of these linear representations.

\paragraph{FP.40 (Proved: odd relative homology over a forgetful base stays forgetful).}
There is a structural obstruction to adding a characteristic-three
Fox module over the forgetful base $Q_{12}$, and it holds for
every finite forgetful base and every odd relative characteristic.
Let $q_L:L=B_4/\langle z_4\rangle\to\widehat B$ be the
fold followed by quotienting the three-strand centre. Then
\[
                  C=\ker(q_L)\simeq C_2\amalg C_2\amalg C_2.
 \tag{FP.40.1}
\]
Indeed the Ferrari kernel is $K=\widehat F(\xi_1,\xi_2)$
by (FP.26). The element $j=(abc)^2$ has order two in $L$,
and its conjugation inverts each $\xi_i$, by squaring
(FP.26.6). The kernel $C$ is $K\rtimes\langle j\rangle$:
upstairs $q((abc)^2)=z_3$, and quotienting $z_4$ makes
this remaining central quotient have order two. Its presentation
\[
 \langle\xi_1,\xi_2,j\mid j^2=1,
           j\xi_1j=\xi_1^{-1},\ j\xi_2j=\xi_2^{-1}\rangle
\]
becomes the free product of three involutions
$j,\xi_1j,\xi_2j$, proving (FP.40.1), discretely and after
completion.

Let $N_B\triangleleft_o\widehat B$ be any finite base kernel
and $N_L=q_L^{-1}(N_B)$. For any odd prime $p$, the three
involution generators of $C$ belong to $N_L^p$, since an
involution equals its own $p$th power. Hence
$C\subseteq N_L'N_L^p$. The natural quotient is therefore
exactly
\[
 \boxed{\quad L/(N_L'N_L^p)
           \simeq\widehat B/(N_B'N_B^p).\quad}
 \tag{FP.40.2}
\]
To see equality of the indicated kernels, quotient $N_L$ by
$C$: commutators and $p$th powers map onto those of $N_B$,
and $C$ is already killed. Thus the purported relative homology
extension contains no additional Brunnian information. Its
induced five-point quotient factors through the five forgetful
maps, where (FP.30) makes the pentagon identically invisible.
The same proof applies to relative exponent $p^r$ for every
$r\geq1$. In particular the proposed characteristic-three
Fox refinement over the order-$576$ congruence base is excluded
without an unnecessary $2304\times1728$ matrix computation.

In fact the obstruction holds for the entire relative pro-$p$
completion. Write $O^p(N)$ for the intersection of all open
normal subgroups of $N$ with $p$-group quotient. Every homomorphism
from $N_L$ to a finite odd $p$-group kills the three involutions,
and hence kills $C$. The universal property gives
\[
 L/O^p(N_L)\simeq\widehat B/O^p(N_B).
 \tag{FP.40.3}
\]
Thus iterating odd relative homology over the same forgetful
base cannot evade this obstruction.

\paragraph{FP.41 (Proved: the first mixed two--three Ferrari cover adds no pentagon equation).}
The preceding obstruction motivates retaining $C^{\rm ab}$
before taking homology in characteristic three. Put
\[
 C=\langle t_0,t_1,t_2\mid t_0^2=t_1^2=t_2^2=1\rangle,
 \qquad
 \overline C=C/(C''(C')^3).
\]
The subgroup $C'$ is free of rank five and index eight, so
$|\overline C|=8\cdot3^5=1944$. In the notation of (FP.40),
$t_0=j$, $t_1=\xi_1j$, $t_2=\xi_2j$. The exact conjugations
by the generators of $L$ are
\[
\begin{split}
 a:(t_0,t_1,t_2)&\longmapsto(t_0t_1t_0,t_0,t_2),\\
 b:(t_0,t_1,t_2)&\longmapsto(t_2,t_1,t_2t_0t_2),\\
 c&=\operatorname{Inn}(t_1t_0)\,a.
\end{split}
 \tag{FP.41.1}
\]
These follow directly from the Ferrari Nielsen formulas in
(FP.26). In particular this constructs an actual homomorphism
$L\to\operatorname{Aut}(\overline C)$.

The character decomposition of $V=C'^{\rm ab}/3$ under
$C^{\rm ab}=(C_2)^3$ is
\[
 V=B\oplus P,
 \qquad
 B=\chi_{110}\oplus\chi_{101}\oplus\chi_{011},
 \qquad P=\chi_{111}^{\oplus2}.
 \tag{FP.41.2}
\]
For completeness, the Cayley graph of the three involutions
has edge module $\bigoplus_i\operatorname{Ind}_{\langle t_i\rangle}^{(C_2)^3}
\operatorname{sgn}$ and vertex module the regular representation.
A nontrivial character of weight $k$ therefore occurs in its
first homology with multiplicity $k-1$. This proves (FP.41.2).
The summand $P$ identifies canonically with the homology of the
index-two Ferrari subgroup $K$, by averaging over the even
subgroup $(C_2)^2$. The other three lines are its boundary
homology after passing to the index-four cover of $K$.

There is a precise collapse in the pentagon test. In a faithful
affine representation of the automorphism group on $V$, all
five $X_i$ act on $B$ by translations, with vectors
\[
 (2,0,0),\quad(0,1,0),\quad(1,0,0),\quad
 (2,1,1),\quad(1,1,2).
 \tag{FP.41.3}
\]
They preserve the direct sum with the affine $P$ block. Their
$P$ blocks, after one common conjugation by translation $(1,1)$,
are precisely the images of the same five words under
\[
 a\longmapsto(0,U),\qquad
 b\longmapsto(0,L_0),\qquad
 c\longmapsto(e_1,U)
 \quad\hbox{in }\mathbb F_3^2\rtimes\operatorname{SL}_2(\mathbb F_3).
 \tag{FP.41.4}
\]
This is exactly the quotient $\eta_3$ of (FP.31), with its
marking preserved, not merely an abstract group isomorphism.

Every $f\in\widehat F'$ is one on an abelian target.
Consequently all five $f$-factors have identity $B$ block,
and their entire pentagon product is exactly the old
$\eta_3$ product. For actual gentle elements the necessary
equation is still $a_f=\rho$ from (FP.33). Adjoining the
forgetful $q_{12}$ component changes nothing: there the
five $X_i$ have images $(x,y,x,1,1)$, and the pentagon
cancels by the two-cycle relation.

More strongly, every evaluation-pair group in this joint
quotient has derived length at most three. The boundary
projection is abelian; the compact pair series are
$24,8,2,1$ for the two harmonic pairs and
$18,3,1$ or $6,3,1$ for the others. The forgetful level-twelve
pair groups also have derived length at most three.
The actual element from LI.8, with $f\in F'''$ and
$\Delta=1$, therefore kills every one of the five factors,
not just their product. The induced five-point quotient has
the same obstruction.

The construction and every automorphism identity on all $1944$
elements are independently checked in the lifting audit's
\texttt{check\_mixed\_ferrari\_quotient.py}. The exact common
conjugacy, all boundary blocks, and all compact pair series
are checked in \texttt{check\_mixed\_ferrari\_collapse.py} and
its output. This result leaves open the next nonabelian
stage, where commutators between $B$ and $P$ are retained.

\paragraph{FP.42 (Proved: the complete class-two Ferrari stage is also insufficient).}
Retain the mixed commutators suggested by (FP.41). Over
$\mathbb F_3$ put $V=B\oplus P$, $W=B\otimes P$ and
\[
 N=V\oplus W,\qquad
 (v,z)(v',z')=(v+v',z+z'+\tfrac12[v,v']),
 \quad
 [v,v']=v_B\otimes v'_P-v'_B\otimes v_P.
 \tag{FP.42.1}
\]
The group $A=(C_2)^3$ acts with the characters in (FP.41.2)
on $V$ and their products on $W$. The six characters on $W$
are the three characters of weight one, each twice. Let
$E=N\rtimes A$. It has order $8\cdot3^{11}=1417176$.
In the coordinates of the certificate, the three elements
\[
 t_0=((0,0),e_0),\quad
 t_1=(((1,0,0,1,0),0),e_1),\quad
 t_2=(((0,1,1,0,1),0),e_2)
\]
are involutions. Their images generate $E/[N,N]$ by (FP.41),
and hence generate $E$: their intersection with $N$ generates
$N/\Phi(N)$, so the finite-$3$-group basis theorem applies.
Thus $E$ is an actual quotient of $C$, not just a module
with a putative action. It is obtained from the full free
class-two exponent-three quotient of $C'$ by killing
$\bigwedge^2B\oplus\bigwedge^2P$.

The two sets of characters on $V$ and $W$ are disjoint.
Consequently $\operatorname{Hom}_A(V,W)=0$, and
$C_N(A)=1$. Schur--Zassenhaus, followed by this vanishing,
shows that every induced automorphism has the unique form
$\operatorname{Inn}(n)$ followed by its graded linear map.
This reduces the exact computation to six additional central
coordinates. The images (FP.41.1) determine those coordinates
uniquely, and give a homomorphism
$L\to N\rtimes\operatorname{GL}(V)$. Every braid,
distant-commutation and full-twist relator is checked in the
actual class-two multiplication, including all nonlinear terms.

There is an explicit structural explanation of the pair groups.
Let $h\in N$ have
\[
 h_V=(0,0,0,1,1),\qquad
 h_W=(0,2,2,0,2,0).
\]
Conjugation by $h$ sends the harmonic pair to
\[
 hX_1h^{-1}=(2e_1,U^2),\qquad
 hX_2h^{-1}=(e_2,L_0^2),
 \tag{FP.42.2}
\]
with zero compact and mixed translations. Hence that pair
embeds into the direct product of an abelian boundary group
and $\operatorname{SL}_2(\mathbb F_3)$. The reversing half
twist gives the other harmonic pair. In each remaining pair,
the two compact linear parts commute. Its first derived group
therefore lies in $P\oplus W$, an abelian subgroup of $N$.
The exact orders and derived series are
\[
\begin{array}{c|c}
 (X_1,X_2),(X_2,X_3)&72,8,2,1\\
 (X_3,X_4),(X_5,X_1),(X_4,X_5)&486,27,1.
\end{array}
 \tag{FP.42.3}
\]
Schreier rewriting and an independent enumeration of each
complete pair group agree. The harmonic pairs have no kernel
in the six new mixed coordinates; each other pair has a
two-dimensional such kernel. Thus the available mixed module
does not produce a coprime action in the harmonic second
derived subgroup.

The same conclusion holds for the \emph{full} class-two
quotient
\[
 E_{\rm full}=C/(\gamma_3(C')(C')^3),
       \qquad |E_{\rm full}|=8\cdot3^{15}.
 \tag{FP.42.4}
\]
Here all powers and subgroups are understood topologically.
The omitted kernel inside the nilpotent subgroup is
$Z_0=\bigwedge^2B\oplus\bigwedge^2P$, of dimension four.
Consider the kernel $\mathcal K$ of the induced map on the
automorphism groups preserving these canonical character
summands. An automorphism in $\mathcal K$ is the identity on
$V$, hence on $\bigwedge^2V$, and is described by a
$Z_0$-valued cocycle. These cocycles add under composition,
so $\mathcal K$ is elementary abelian of characteristic three.
Its two types of coordinates are
\[
 \operatorname{Inn}(\bigwedge^2B),\qquad
 \operatorname{Hom}_A(B,\bigwedge^2B),
 \tag{FP.42.5}
\]
each of dimension three. The invariant line
$\bigwedge^2P$ is central in $E_{\rm full}$ and gives no
additional inner automorphism. There are no other equivariant
maps from $V$ to $Z_0$ by their character decompositions.

One must distinguish this automorphism kernel from the
four-dimensional central kernel in the nilpotent subgroup.
A pure $X_i$ can shear the derivation coordinate in (FP.42.5)
by an inner correction, because its boundary translation need
not vanish. Every element of the second derived pair group,
however, has zero boundary translation. Its boundary linear
part is the identity, and its compact part is $\pm I$ in the
harmonic case, or the identity in the other cases.
It therefore acts trivially on both summands in (FP.42.5).
Explicitly, write a kernel element as
$\operatorname{Inn}(z)\,d_D$, where
$d_D(v,w)=(v,w+D(v_B))$. Conjugation by an element with
boundary linear part $I$ and boundary translation $u_B$
sends $(z,D)$ to $(z-D(u_B),D)$; the compact scalar $\pm I$
does not affect these coordinates. This proves the asserted
triviality when $u_B=0$.
The second derived subgroup in the full quotient is a central
extension of $C_2$ (harmonic case), or of the trivial group,
by a subgroup of $\mathcal K$. A central extension of a cyclic
group is abelian. Consequently all five full pair images still
have derived length at most three.

The actual LI.8 element with $f\in F'''$ and $\Delta=1$
therefore satisfies every pentagon factor separately in both
class-two quotients, also after adjoining any of the previous
forgetful level-twelve factors. No assertion that this element
belongs to GT is made. The exact nonlinear construction is
\texttt{agent\_lifting/code/check\_mixed\_ferrari\_class2.py};
the independent pair proof and complete enumeration are
\texttt{check\_mixed\_class2\_pairs.py} and its output.

\paragraph{FP.43 (Proved, with LI.8): the entire relative pro-three Ferrari completion is blind to an actual mixed-parity witness.}
The class-two obstruction has a uniform structural extension.
Put
\[
 N=C'/O^3(C'),\qquad
 C_\infty=C/O^3(C'),\qquad
 L_\infty=L/O^3(C').
 \tag{FP.43.1}
\]
Then $N$ is free pro-three of rank five, and
$C_\infty/N=(C_2)^3$. The kernel $O^3(C')$ is characteristic
in $C$ and normal in $L$, so these are actual profinite
quotients with the given braid action.

The restriction
$\operatorname{Aut}(C_\infty)\to\operatorname{Aut}(N)$
is injective. Indeed $C_{C_\infty}(N)=1$: its intersection
with $N$ is the centre of a nonabelian free pro-three group,
and every nonzero element of $(C_2)^3$ acts nontrivially on
$N^{\rm ab}/3$ by the five characters in (FP.41.2).
An automorphism fixing $N$ pointwise therefore fixes every
element of $C_\infty$. The kernel of
$\operatorname{Aut}(N)\to\operatorname{GL}(N^{\rm ab}/3)$
is pro-three, by the finite-$p$-group Frattini automorphism
theorem and inverse limits. In particular $C_\infty$ itself
is centreless.

Let $G$ be the closed image of $\langle a,b\rangle$ in
$\operatorname{Aut}(C_\infty)$. Its linear image is exactly
\[
 G_0=S_3\times\operatorname{SL}_2(\mathbb F_3),
 \qquad |G_0|=144,\quad |G_0'|=24,\quad G_0''=C_2.
 \tag{FP.43.2}
\]
This is a marked assertion: the $S_3$ is the boundary action,
the second factor is the compact action, and
$z_3=(ab)^3$ maps to $(1,-I)$.
One can also see the direct product without enumeration:
the boundary images of $a,b$ have order two whereas their
compact images have order three, so their third and fourth
powers separate the two factors.

Write $J$ for the pro-three kernel of $G\to G_0$.
The element $z_3$ is central in $G$. Its closed procyclic
subgroup has a unique two-primary component, an involution
$\iota$ mapping to $(1,-I)$. Consequently
\[
      \operatorname{pr}^{-1}(G_0'')
          =J\times\langle\iota\rangle
\]
is pronilpotent. The harmonic pair subgroup
$\overline{\langle X_1,X_2\rangle}$ has its second derived
subgroup in this preimage. The involution need not belong to
the harmonic pair itself; containment in the displayed
pronilpotent ambient group is sufficient. The reversing
half twist handles $(X_2,X_3)$. For each of the other three
pairs the linear image is abelian, since its compact generators
are $U^2$ and $-I$, or two copies of $-I$, and its boundary
linear action is trivial. Its first derived subgroup is
therefore pro-three. Thus all five evaluation-pair second
derived subgroups are pronilpotent, uniformly at every depth.

Choose the actual element $k=(1,f_k)$ from LI.8, for which
\[
 f_k\in R_{\rm nil}(F''),\qquad D(k)=1,
                   \qquad\Delta(k)=1.
 \tag{FP.43.3}
\]
Every continuous evaluation $F''\to P''$ into one of these
pair groups kills $R_{\rm nil}(F'')$. Hence each of the five
$f_k$-factors acts trivially on $C_\infty$.

There is a stronger conclusion in the actual quotient
$L_\infty$, retaining the full forgetful base $\widehat B$.
The pentagon product has trivial image in $\widehat B$ by
(FP.30), so belongs to $C_\infty$. Its conjugation action on
$C_\infty$ is trivial by the preceding paragraph. Centrelessness
then implies that this product is one in $L_\infty$.
The cyclically induced five-point quotient has the same
property. Individual factors are asserted to be one in the
\emph{action} quotient, and in its joint level-twelve target;
in $L_\infty$ the assertion is about their pentagon product.

This excludes the entire proposed relative pro-three Ferrari
family as a pentagon certificate for $\Delta$, not merely its
first few nilpotent truncations. It does not assert that $k$
satisfies the pentagon in the full profinite five-point group.
The complete independent proof is recorded as FR.1--FR.4 in
\texttt{agent\_literature/Ferrari\_relative\_pro3.tex}; the
lifting audit independently checked the pro-three and central
involution steps. Exact finite linear data are included in
\texttt{check\_mixed\_ferrari\_collapse.py} and its output.

\paragraph{FP.44 (Proved: an actual rank-129 mixed-prime homology quotient).}
Let $C=C_2\amalg C_2\amalg C_2$ as above and put
\[
 D=C''(C')^2,\qquad Q=C/D,
 \qquad V_D=D^{\rm ab}/3.
 \tag{FP.44.1}
\]
Since $C'$ is free of rank five, $C'/D=(C_2)^5$ and
$|Q|=8\cdot32=256$. The subgroup $D$ is free of rank
$1+32(5-1)=129$. All these subgroups are characteristic
in $C$ and hence invariant under the literal Ferrari action
of $L$. Conjugation therefore defines the actual continuous
homomorphism
\[
           \eta_D:L\longrightarrow
                    \operatorname{GL}(V_D)
                    =\operatorname{GL}_{129}(\mathbb F_3).
 \tag{FP.44.2}
\]
This target has an additional two-primary layer inside $C'$
before its relative three-primary step; it is not covered by
the relative-pronilpotent obstruction for $C'$ itself.

Here is an exact finite construction. The three-cube has
eight vertices and twelve unoriented edges. Over $\mathbb F_2$,
map the three involutions to $(e_i,t_i)$ in its edge-chain
module semidirect $(C_2)^3$. Their image consists of the
$256$ path chains with boundary equal to their endpoint minus
the basepoint. This realizes $Q$ with its prescribed marking.
Apply the exact involution-word automorphisms (FP.41.1) to
obtain the three automorphisms of $Q$.

The Schreier graph of $Q$ with its three involution generators
has $256$ vertices and $384$ unoriented edges. A spanning
tree leaves $129$ chords. Their based cycles give a literal
basis of $V_D$. For each automorphism, replace every oriented
edge by the path spelling the image of its involution, and
read the chord coordinates. This gives the matrices of
$a,b,c$, without a guessed representation or an unproved
descent. The certificate verifies every vertex action, every
edge endpoint, every cycle boundary, and the full intertwining
of edge chains and chord coordinates. It then verifies all
three braid/commutation relations and $(abc)^4=1$ on
$V_D$. The five $X_i$ have exact matrix orders
\[
                         6,6,6,4,4.
 \tag{FP.44.3}
\]
The construction is \texttt{check\_ferrari\_rank129.py}; its
output and the complete matrices and graph are retained in
\texttt{check\_ferrari\_rank129.out} and
\texttt{ferrari\_rank129.npz}.

\paragraph{FP.45 (Proved: the harmonic second derived image is not nilpotent).}
In the actual quotient (FP.44), set $x=\eta_D(X_1)$,
$y=\eta_D(X_2)$ and
\[
 c=[x,y],\qquad w=[c,xcx^{-1}],\qquad
 u=w^6,\qquad v=(xy^{-1})u(xy^{-1})^{-1}.
 \tag{FP.45.1}
\]
The exact matrix orders are
\[
                  |w|=12,\qquad |u|=|v|=2,
                              \qquad |uv|=6.
 \tag{FP.45.2}
\]
Both $u,v$ belong to the second derived subgroup of
$\langle x,y\rangle$, because $w$ is a commutator of two
elements of its first derived subgroup. In a finite nilpotent
group the two-elements form the unique Sylow two-subgroup;
the product of two involutions must therefore have two-power
order. Equation (FP.45.2) is an exact contradiction.
Thus the pair-second-derived nilpotence test that excluded
all preceding Ferrari truncations fails for this target.

The source words, orders and matrix products are checked in
\path{check_rank129_pair_gate.py}. An independent audit
uses the complete free word, integer-64 arithmetic and the
known identities $x^{-1}=x^5$, $y^{-1}=y^5$, with no imported
inverse routine or order-search implementation. Its source
is \texttt{agent\_lifting/code/check\_rank129\_gate\_audit.py}.
This proves a property of an actual finite quotient. It is
neither a nonzero value of $\Delta$ on GT nor evidence that
an arbitrary finite parameter assignment lifts.

\paragraph{FP.46 (Proved: the obstruction is localized in two-branch boundary/proper coupling).}
The deck group $C'/D=(C_2)^5$ acts semisimply on $V_D$.
The trivial character has multiplicity five, and each of its
$31$ nontrivial characters has multiplicity four. The exact
simultaneous eigenspaces are computed from deck translations
on the same Schreier graph, with no new homological marking.

There are $64$ cusp loops in the $D$ cover. Writing
$w_0=t_1t_0t_2$, they are the cycles $q w_0^4q^{-1}$,
one for each of the $64$ cosets modulo the order-four image
of $w_0$ in $Q$. Their span has rank $63$ and is invariant
under $a,b,c$. The quotient, the proper homology, has rank
$66$. On each nontrivial deck character the possible
dimensions are
\[
\begin{array}{c|c|c|c}
 \text{branch points among the four cusps}&
 \text{number of characters}&\text{boundary}&\text{proper}\\\hline
 0&3&4&0\\
 2&24&2&2\\
 4&4&0&4.
\end{array}
 \tag{FP.46.1}
\]
These counts are verified by evaluation on the four actual
cusp words; equivalently they follow from Riemann--Hurwitz
for the corresponding double covers. The full $L$ character
orbits have sizes $24,3,3,1$ among the nontrivial characters:
the final singleton is one of the fully branched covers.

The involution witness (FP.45) is supported entirely on the
$24$ two-branch characters. It vanishes on both fully branched
pieces and on the unramified pieces. On the proper rank-$66$
quotient the product of its two involutions has order two;
on the open homology it has order six. Thus the new
non-nilpotence uses the extension of proper homology by
boundary homology and is lost by passing to the compact
Jacobian representation alone.

In the exact binary character basis retained by the code, the
two harmonic orbits
\[
          (1,7,18,20),\qquad (8,14,27,29)
 \tag{FP.46.2}
\]
each have dimension $16$, with boundary and proper dimensions
eight each. On either orbit the two involutions in (FP.45)
have product of exact order three. The first orbit's literal
harmonic image is an abelian extension with
\[
 |Q_{\rm lin}|=4608,\quad
 |Q_{\rm lin}'|=384,\quad |Q_{\rm lin}''|=64,
 \quad |Q_{\rm lin}'''|=4,
 \qquad U\simeq\mathbb F_3^4,
\]
so its full order is $4608\cdot3^4=373248$.
Here $Q_{\rm lin}$ retains the separate boundary and proper
blocks, while $U$ is the exact image of the off-diagonal
extension module. Schreier loops in the finite base compute
the entire four-dimensional kernel, not a chosen subspace.

The exact decomposition, branch tests, cusp spans and small
module are respectively
\texttt{check\_rank129\_deck\_blocks.py},
\texttt{check\_rank129\_prym.py},
\texttt{check\_rank129\_boundary.py}, and
\texttt{check\_rank129\_harmonic\_module.py}, with their
outputs and complete matrix data retained alongside them.
These marked data identify a new coupling; they do not yet
identify its character with $\Delta$.

\paragraph{FP.47 (Proved necessary matrix identity, subsequently proved automatic on gentle elements).}
Fix character $1$ in (FP.46.2). All five evaluation-pair
permutation images on its character orbit are abelian.
Consequently every $f\in\widehat F'$ preserves this individual
four-dimensional character space in each of the five factors.
For $\Phi=(1,f)\in\mathcal H$, one has the stronger elementary
fact $f\in\gamma_3(\widehat F)$: in the class-two quotient
write $f=[x,y]^e$; the multiplier-one three-cycle gives
$3e=0$ in $\widehat{\mathbb Z}$, hence $e=0$.

The literal boundary pair images have lower-central order
series $8,2,1$ or $4,1$. Thus all five boundary $f$-blocks
are the identity for every actual gentle element. The proper
pair images for $(X_3,X_4)$ and $(X_5,X_1)$ have order $24$
with central derived subgroup of order two, and that for
$(X_4,X_5)$ has order eight with the same property. Their
proper $f$-blocks are also the identity. Write the five
four-dimensional blocks, in pentagon order, as
\[
 \begin{pmatrix}I&C_{34}\\0&I\end{pmatrix},\quad
 \begin{pmatrix}I&C_{51}\\0&I\end{pmatrix},\quad
 \begin{pmatrix}I&C_{23}\\0&P_{23}\end{pmatrix},\quad
 \begin{pmatrix}I&C_{45}\\0&I\end{pmatrix},\quad
 \begin{pmatrix}I&C_{12}\\0&P_{12}\end{pmatrix}.
\]
All entries are over $\mathbb F_3$; each $C_{ij}$ and $P_{ij}$
is $2\times2$ and is the evaluation of the \emph{same} source
word $f$ in its specified pair. The full pentagon implies
the exact equations
\[
 \boxed{\quad P_{23}P_{12}=I,\qquad
 C_{34}+C_{51}+(C_{23}+C_{45})P_{12}+C_{12}=0.\quad}
 \tag{FP.47.1}
\]
This follows by literal block multiplication. The two harmonic
proper pair groups have derived series
$2304,192,64,4,1$; their values must not be replaced by
independent parameters. All the group orders and class-two
assertions are checked in
\texttt{check\_rank129\_pentagon\_blocks.py} and its output.

Equation (FP.47.1) is a literal boundary/proper matrix identity
in an actual five-point representation. The subsequent exact
normalization in (FP.48)--(FP.49) proves that it holds for
every multiplier-one gentle element. It therefore supplies
no new comparison between $\kappa_d$ and $\epsilon_2(W)$.
The separate affine extension in (FP.50) retains information
that this homology action kills. No implication from a
surviving finite matrix assignment to a genuine GT lift is
asserted.

\paragraph{FP.48 (Proved: the exact open-Prym normalization and its automatic pentagon).}
Each of the $24$ two-branch four-dimensional character
spaces in (FP.46) has a canonical splitting of dimensions
$1+3$. Exactly one of the three deck involutions $t_i$
stabilizes that character with eigenspace dimensions one
and three. The smaller eigenspace is a boundary line; the
larger contains one boundary line and the proper
two-dimensional homology. Computing these deck involutions
on the original graph proves that every $a,b,c$ transition
preserves the two summands, including their character
permutations. Thus the rank-$96$ piece splits as a
rank-$24$ boundary representation plus a rank-$72$ open-Prym
representation.

In the latter, each character transition has a $3\times3$
block $\bigl(\begin{smallmatrix}\alpha&v\\0&M\end{smallmatrix}\bigr)$,
with $\alpha=\pm1$. Tensoring by the inverse boundary line
divides the entire block by $\alpha$. This is a genuine
representation operation on the character bundle, not an
independent rescaling of selected generator matrices. A
single consistent change of the proper basis on each of the
$24$ characters then makes every proper transition have
determinant one. Its existence and consistency around all
generator loops are checked exactly. The resulting
homomorphism is a block-permutation representation with
blocks in $\operatorname{ASL}_2(\mathbb F_3)$, in the
row-vector convention
\[
                  \begin{pmatrix}1&v\\0&M\end{pmatrix},
                         \qquad \det M=1.
\]
Its full $L$ permutation image has order $192$; the five
pure $X_i$ have an elementary abelian permutation image of
order eight. All defining relations of $L$ hold after the
tensor operation and the fixed basis change.

For the untwisted boundary line on \emph{all} $24$ characters,
the five pair groups have lower-central order series, in
pentagon order,
\[
        (16,2,1),\ (16,2,1),\ (8,2,1),\
        (16,2,1),\ (8,2,1).
\]
Since every actual gentle $f$ belongs to $\gamma_3F$ by
(FP.47), all its five boundary-line evaluations are one.
The tensor normalization therefore preserves those actual
$f$-values. This preservation is not claimed for arbitrary
commutator words.

The normalized matrices satisfy the exact marked identities
\[
 X_3=X_1,\qquad X_5=X_4,\qquad
 [X_4,X_1]=[X_4,X_2]=1.
 \tag{FP.48.1}
\]
They also follow from source identities. Put
$u=\xi_1$, $v=\xi_2$, $s=aba$. Then
$c=ua$, $a(u)=u$, $s(u)=v^{-1}$ and $s(v)=vuv^{-1}$, whence
\[
 X_3X_1^{-1}=u^2\in C',\qquad
 cbc=uv^{-1}s,\qquad
 X_5X_4^{-1}=uv^{-1}u^{-1}v^{-1}\in C'.
\]
Every $C'$ deck element acts by the same scalar on an
entire character block and its boundary line, so the tensor
operation kills $C'$ itself. Also $X_4$ is the literal
three-strand centre and commutes with $X_1,X_2$.
Equation (FP.48.1) makes the pentagon cancel by commuting
arguments and $f(y,x)=f(x,y)^{-1}$. Thus the apparently new
row-vector equation from the open-Prym coupling is automatic
on all actual gentle elements.

The exact involution eigenspaces and invariant splitting are
\texttt{check\_rank129\_open\_prym.py}; the determinant gauge,
all $24$ boundary-line tests and the marked identities are
\texttt{check\_open\_prym\_affine.py}. Complete matrices are
\texttt{ferrari\_open\_prym72.npz} and
\texttt{ferrari\_open\_prym\_affine72.npz}.

\paragraph{FP.49 (Proved: the entire linear rank-129 pentagon is automatic).}
The preceding argument also settles the other character
pieces without a separate parameter enumeration. In $L/C'$,
all $X_i$ centralize $C/C'$, while $X_4,X_5$ lie in this
abelian normal subgroup and $X_3$ differs from $X_1$ by
one of its elements. Hence the pentagon of every commutator
two-cycle word lies in $C'$.

For an actual multiplier-one gentle word, (FP.48) shows
that its pentagon has scalar value one on all $24$ two-branch
characters of $C'/D$. These characters span the full
five-dimensional dual of $C'/D$; this is checked in the exact
binary character basis. Therefore
\[
        \boxed{\quad P(f)\in D=C''(C')^2
                   \quad\text{for every }(1,f)\in\mathcal H.\quad}
 \tag{FP.49.1}
\]
Conjugation by an element of $D$ is trivial on $D^{\rm ab}$.
Consequently the pentagon is one in the rank-$129$ homology
\emph{action}, over every coefficient ring, including
$\widehat{\mathbb Z}$. This establishes universal blindness
of that linear test despite the non-nilpotent pair subgroup
proved in (FP.45). The intrinsic membership statement
(FP.49.1) must be distinguished from membership in
$D'D^3$, which is the next genuine affine test.

\paragraph{FP.50 (Proved affine construction and nontrivial Brunnian check; the actual-lift obstruction is proved in FP.52).}
The exact graph construction also gives the faithful
relative translation coordinates omitted by its homology
action. Let $C_1$ be the $384$-dimensional edge-chain space
of (FP.44), $M=V_D$ its cycle space, and $p_q$ the chosen
tree path from the identity to $q\in Q$. Over $\mathbb F_3$ put
\[
 u=-\frac1{256}\sum_{q\in Q}p_q,
 \qquad \partial u=e_1-\frac1{256}\sum_{q\in Q}e_q.
 \tag{FP.50.1}
\]
The Fox path-chain model, conjugated by $u$, identifies
\[
          C/(D'D^3)\simeq M\rtimes Q.
\]
For a literal graph chain automorphism $T_g$, define
\[
                 r_g=u-T_gu\in M,
                 \qquad M_g=T_g|_M.
 \tag{FP.50.2}
\]
The boundary of $r_g$ is zero because $g$ fixes the identity
vertex and permutes all vertices. The exact cocycle identity
$r_{gh}=r_g+M_gr_h$ proves that
\[
             g\longmapsto(r_g,M_g)
       \quad\hbox{defines } L\longrightarrow
             \mathbb F_3^{129}\rtimes\operatorname{GL}_{129}(\mathbb F_3).
 \tag{FP.50.3}
\]
All braid and full-twist relations are checked in this
affine group, including the translations. The five $X_i$
have orders $6,6,6,12,12$, as opposed to the linear orders
in (FP.44).

For an inner element $d\in D$ with homology class $m\in M$,
one has $M_d=I$ and $r_d=m$. Indeed $M^Q=0$, since averaging
identifies the invariant homology with $H_1(C,\mathbb F_3)=0$;
the three generators of $C$ are involutions. Averaging the
exact inner path formula gives $T_du=u-m$.
Combining this with (FP.49), the affine pentagon is therefore
precisely the class
\[
                    [P(f)]\in D/(D'D^3).
 \tag{FP.50.4}
\]
It vanishes if and only if $P(f)\in D'D^3$.
The full forgetful base can be retained in a joint quotient.

This affine target has a concrete nonzero Brunnian check:
$[[X_1,X_2],[X_3,X_4]]$ has identity linear part and a
nonzero translation. The full vector is retained in the
certificate. There is also a word in $F''$ satisfying the
two-cycle identity whose pentagon has a nonzero translation;
that test word is explicitly \emph{not} asserted to satisfy
the gentle three-cycle or define a GT element. These checks
show that the affine construction retains actual source
information killed by the linear representation. They do
not yet produce a gentle obstruction character, a genuine
GT counterexample, or a proof that the affine equation
forces $\Delta=0$.

The chain maps, gauge and affine generators are saved in
\texttt{ferrari\_rank129\_affine.npz}; the construction and
all affine relator checks are part of
\texttt{check\_ferrari\_rank129.py}. The exact Brunnian and
two-cycle probes are
\texttt{check\_rank129\_affine\_defect.py} and its output.
The additional literal-word certificate
\texttt{check\_rank129\_inner\_translation.py} takes the
first chord generator of $D$, substitutes its full conjugated
involution words on every graph edge, and verifies identity
linear action and translation exactly equal to the first
homology basis vector. This independently excludes a mere
coboundary of the $129$-dimensional linear representation.
The structural quotient and faithful-shear proof are also
recorded as FA.1--FA.3 in the literature audit's
\texttt{Ferrari\_affine\_D.tex}.

\paragraph{FP.51 (Proved: no common affine fixed point, even after restriction to commutators).}
The translations in (FP.50) cannot be eliminated by a common
change of origin. The exact systems
$(I-M_g)h=r_g$ for the indicated two generators have the
following ranks over $\mathbb F_3$:
\[
\begin{array}{c|c|c}
 \text{generators}&\text{coefficient rank}&\text{augmented rank}\\\hline
 a,b&125&126\\
 X_1,X_2&111&112\\
 X_2,X_3&111&112\\
 X_3,X_4&117&118\\
 X_4,X_5&120&121.
\end{array}
\]
There is also no common fixed point for any of these pairs'
derived subgroups. For the harmonic pairs, the equations for
the basic commutator and a finite set of its generator
conjugates already have ranks $90,91$; for the last two
pairs they have ranks $64,65$.
These are ranks of explicit inconsistent subsystems, not
claims about the rank of an unenumerated infinite system.
The certificate transforms an equation $Eh=b$ under an
affine generator $(v,M)$ to
$EM^{-1}h=b+EM^{-1}v$ and closes its row space under both
generators and their inverses until contradiction.
Thus even a formula $r_{f(g,h)}=(I-M_{f(g,h)})h_0$ valid
for every commutator word cannot remove the new translations.
This finite cohomological obstruction is compatible with the
actual mixed-parity witness in (FP.52); it does not force
$\Delta$ to vanish. Exact systems and solutions or rank
contradictions are in
\texttt{check\_rank129\_fixed\_points.py} and its output.

\paragraph{FP.52 (Proved with R2: the actual affine rank-129 test also admits a mixed-parity gentle witness).}
The affine test is not killed by the previous pairwise
nilpotence theorem, but a controlled source seed gives a
stronger obstruction. Let $e_3\in\widehat{\mathbb Z}$ have
three-adic component one and every other component zero.
For $z=(xy)^{-1}$ define the genuine profinite words
\[
 A(x,y)=[[x^{e_3},y^{e_3}],[y^{e_3},z^{e_3}]],\qquad
 S=A(y,z)A(x,y)A(z,x).
 \tag{FP.52.1}
\]
Both belong to $F''$. The exponent is $4$ modulo twelve;
the exact signed calculation gives
\[
                            q_{12}(S)=5I.
\]
The seed lies in every section-47 core used in LI.8:
in each peripheral marking one boundary maps to the
order-two generator, whose $e_3$th power is one, killing
the outer double commutator.

Evaluate this \emph{same} word in all five actual affine
pair images and all six peripheral markings. Each image
$s$ has order dividing six and satisfies the following
exact conditions:
\[
 \begin{split}
 s^3&\text{ centralizes both pair generators},\\
 s^4&\text{ acts as the identity on the boundary rank-$63$
        and proper rank-$66$ graded pieces}.
 \end{split}
 \tag{FP.52.2}
\]
The images are one in the three nonharmonic pairs and have
order six in the harmonic pairs. The affine translations
are retained throughout this check.

The conditions have an exact normal-closure consequence.
Inside the affine group preserving the boundary subspace,
the kernel of the map to the boundary and proper graded
linear groups is a finite three-group: in the augmented
matrix representation it consists of unipotent upper block
triangular matrices. It is normal in every pair image.
Thus all conjugates of $s^4$ remain in one normal three-group,
whereas $s^3$ is a central two-element. Since
$s=s^3s^4$, the normal closure of all six peripheral seed
images is a product of a central two-group and a three-group,
and is nilpotent. The six markings account for the
$B/F\simeq S_3$ outer action; conjugations by $F$ are already
conjugations inside the pair. Hence this proves the needed
statement for the full $\widehat B$-normal source closure,
not merely for separate word values.

For the finite joint evaluation map, take the inverse image
of this nilpotent normal closure and intersect with any
open $\widehat B$-normal subgroup containing $F''$.
The resulting open normal base still contains $S$, retains
$q_{12}(S)=5I$, and has nilpotent image in each affine
evaluation target. Applying R2 exactly as in LI.8 and taking
the same compactness limit gives an actual element
$k=(1,f_k)\in\mathcal H$ with
\[
 \Delta(k)=1,\qquad D(k)=1,\qquad
 f_k,\beta_k\in R_{\rm nil}(F''),
\]
whose five affine $f_k$-factors are all one. Since
$P(f_k)\in D$ by (FP.49) and the affine translation is its
faithful class in $D/(D'D^3)$ by (FP.50),
\[
                       P(f_k)=1\quad\text{in }L/(D'D^3).
 \tag{FP.52.3}
\]
The induced five-point quotient has the same obstruction.
This is an actual global gentle witness, obtained by the
accepted lifting theorem from an explicit source seed. It
is not asserted to be a GT element and does not resolve
the full discrepancy question.

The original exact checks are
\texttt{agent\_lifting/code/check\_rank129\_affine\_seed.py}.
Independent arithmetic verifies all $30$ pair/core cases,
the primary decomposition, centrality, full graded blocks
and the signed level-twelve residue in
\texttt{check\_rank129\_affine\_escape\_audit.py} and its
output. The normal-closure argument, rather than survival
of a finite parameter assignment, is the input to R2.

\paragraph{FP.53 (Proved with R2: one witness survives the full relative abelianization of $D$).}
The conclusion of (FP.52) extends to one actual gentle element with
\[
 \Delta(k)=1,\qquad D(k)=1,\qquad P(f_k)=1\quad\hbox{in }L/D'.
 \tag{FP.53.1}
\]
Here $D=C''(C')^2$ remains fixed. This concerns its full abelianization,
not a nonabelian pro-three completion. The compactness proof is
RA.5--RA.6 in
\texttt{agent\_lifting/rank129\_collapse\_affine\_audit.tex}.
The additional inputs have the following independent integral audit.

On $M=H_1(D,\mathbb Z_3)$ the deck average
$E_0=\frac1{32}\sum_{g\in C'/D}g$ is an $L$-equivariant idempotent
of rank five. Exact graph pushforward gives
\[
 JM_g=M_g^{(8)}J,\qquad
 Jr_g=r_g^{(8)}+(I-M_g^{(8)})\delta\quad(g=a,b,c),
 \tag{FP.53.2}
\]
with a saved vector $\delta$ of denominator $256$. Thus this affine
summand is conjugate by one fixed translation to the old rank-five
affine action. In all thirty pair/core cases, $S^3$ is the identity
on the complementary rank-$124$ affine summand modulo three,
including translations. The congruence kernel is pro-three; hence
the whole complementary cyclic image of $S$ is pro-three.
The old rank-five factor has the central braid involution from
(FP.43). The three-primary seed parts lie in the normal preimage
of the boundary/proper unipotent kernel. These facts establish the
pronilpotent normal-image condition at every three-adic level.

The exact integral augmented matrices satisfy
$(X_i^2-I)^2=0$ for $i\leq3$ and $(X_i^4-I)^2=0$ for $i=4,5$;
the third peripheral matrices likewise have a power two or four
with square-zero difference from the identity. Away from two and
three their procyclic orders have primes only in $\{2,p\}$, so
their $e_3$ powers vanish. At two no division by $256$ is used:
$C'/(D'D^{2^n})$ is a finite two-group, and the retained condition
modulo $O^2(C')$ implies all these conditions. R2 and compactness
give (FP.53.1). Certificates:
\texttt{build\_rank129\_integral.py},
\texttt{check\_rank129\_integral\_peripherals.py},
\texttt{check\_rank129\_transfer.py}, and the updated
\texttt{check\_rank129\_affine\_escape\_audit.py}.

\paragraph{FP.54 (Proved: an actual full-chain fourth-power cover target).}
Put $D_4=C''(C')^4$. Then $C'/D_4\simeq(\mathbb Z/4)^5$,
$[C:D_4]=8192$, and $D_4$ is free of rank $4097$.
Use the five cube-Schreier generators $q_1,\ldots,q_5$ of $C'$,
and $\mathbb F_9=\mathbb F_3[i]$, $i^2=-1$.
For every row character $\chi\in(\mathbb Z/4)^5$, retain all five
Fox chain coordinates, with boundary
$d_\chi=(i^{\chi_1}-1,\ldots,i^{\chi_5}-1)$.
An automorphism $g$ with integral abelianization $T_g$ sends source
$\chi$ to target $\eta=\chi T_g^{-1}$ with matrix $J_g(\eta)$.
The identities $d_\eta J_g(\eta)=d_\chi$ and the Fox chain rule
give an actual monomial representation of $L$.
All braid, distant, and full-twist relations hold both in the
literal free substitutions and in the evaluated matrices.

For nontrivial $\chi$, the cycle subspace has dimension four.
The fifth chain coordinate is essential: for $h\in D_4$,
\[
                  J_{\operatorname{Ad}(h)}(\chi)
                    =I-v_h(\chi)d_\chi.
 \tag{FP.54.1}
\]
Thus the full chain action detects the corresponding translation,
while restriction to cycles kills it. Adding the earlier rank-five
affine factor detects the trivial character. The faithful source
construction is FF.1--FF.3 in
\path{agent_literature/Ferrari_fourth_power_Fox.tex}.
No unproved stability under a GT action is used.

The pure character actions need not commute. In pentagon order
$34,51,23,45,12$, their pair groups have orders $8,8,32,2,32$.
The three nonharmonic groups are abelian; both harmonic groups
have lower central and derived orders $32,2,1$.
The full-chain $X_i$ have exact orders $12,12,12,24,24$.
The four cusp words in the fixed basis are
\[
 q_1q_4^{-1}q_2^{-1},\quad q_1q_3q_5q_1^{-1},\quad
 q_2q_4q_5^{-1}q_2^{-1},\quad
 q_1q_5^{-1}q_3^{-1}q_2q_1^{-1}q_5q_1^{-1}.
\]
Their Fox cycles give the equivariant boundary/proper filtration.
Its total dimensions are $1023$ and $3074$, adding to $4097$.
Certificates are \texttt{check\_ferrari\_d4\_fox.py},
\path{check_d4_seed_normal_closure.py}, and the literal input
\path{agent_lifting/records/f5_nielsen_fox.json}.

\paragraph{FP.55 (Refuted: the fixed cyclic seed has nilpotent normal image at the fourth-power stage).}
The seed $S$ of (FP.52) does \emph{not} have nilpotent normal image
in (FP.54). This concerns that exact controlled source word;
it neither excludes another R2 seed nor settles the main question.
For $x=X_1$, $y=X_2$, set
\[
              u=S(x,y)^3,\qquad v=(yx)u(yx)^{-1}.
 \tag{FP.55.1}
\]
All peripheral orders divide $24$, so $e_3$ is represented by
$16$ in this literal evaluation. The seed has order six; $u,v$
are involutions fixing all character indices. At
$\chi=(1,0,1,1,0)$ the cusp exponents are $(0,1,1,2)$.
The invariant boundary/cycle/chain filtration has dimensions
$1,4,5$. In its proper three-dimensional quotient the matrices
of $u,v$ over $\mathbb F_3[i]$, $i^2=-1$, are
\[
 U=\begin{pmatrix}-1&1&i\\0&-1+i&-1\\0&-1+i&1-i\end{pmatrix},\qquad
 V=\begin{pmatrix}1-i&-i&-1-i\\i&-i&i\\-1-i&-i&1-i\end{pmatrix}.
 \tag{FP.55.2}
\]
They satisfy $U^2=V^2=I$, and $UV$ has exact order three.
The cycle-space and full-chain products have exact order nine.
The obstruction therefore already survives a proper quotient;
it is not an artifact of the fifth chain direction.

The minimal harmonic character orbit containing this block has
sixteen indices
\[
81,86,91,92,386,391,392,397,755,756,761,766,800,805,810,815,
\]
with index $\sum_j4^{j-1}\chi_j$.
The normal closure of $S(x,y)$ contains both involutions in
(FP.55.1). In a finite nilpotent group all two-elements belong
to its Sylow two-subgroup, and their product remains a two-element.
The proper order-three image contradicts this. The product in
the whole $1024$-block target has order eighteen, giving the same
contradiction without passing to a block subquotient.

\path{check_d4_seed_obstruction.py} independently reconstructs
the same source word and saves the exact five-by-five matrices,
filtration basis, and the three order calculations.
The minimal orbit and proper blocks are in
\texttt{ferrari\_d4\_obstruction\_orbit.npz} and
\texttt{ferrari\_d4\_proper16.npz}.
The bounded conjugacy exploration in the auxiliary normal-closure
script is not used as a proof; the two explicit involutions suffice.

\paragraph{FP.56 (Proved: all active primitive proper images are simple, and the source seed can be corrected).}
Let $N=\langle\!\langle S\rangle\!\rangle_{\widehat B}$ be the
closed normal subgroup generated by the same actual source seed as
in (FP.52). In each harmonic evaluation, each of its six peripheral
seed markings has precisely forty-eight nontrivial primitive proper
blocks, all of dimension three, together with the old proper
two-dimensional block at the trivial character. The complete
transported primitive support consists of forty harmonic character
orbits: thirty-two of size eight and eight of size sixteen, totaling
$384$ characters. No other primitive proper dimension is active.

For each of these forty orbits, the local proper image of $N$ has
order $6048$ and is nonabelian simple. This is an exact assertion
about the simultaneous-source normal subgroup, not merely the
normal closure of one independently chosen matrix. The computation
uses all six peripheral images, transports every character to one
orbit representative, and closes their image under the exact
Schreier stabilizer generators. The stabilizer groups have orders
$6048$, $12096$, or $24192$; in every case the indicated normal
image has order $6048$.

Simplicity is independently certified rather than inferred from
order or perfectness. The reference group has fourteen conjugacy
classes; the normal closure of every nonidentity class is the
whole group. For each of the forty normal images, the certificate
finds its invariant nondegenerate Hermitian form over $\mathbb F_9$
and an explicit orthonormal basis. In these bases every generator
lies in the same reference group. Equal orders give equality,
so all forty groups are explicitly conjugate to the certified
centreless simple group.

The full-chain assertions needed to use this classification are
also checked. All thirty pair/core values fix character indices,
boundary blocks, and augmentation quotients. The three nonharmonic
pair evaluations of $S$ are one. Moreover, conjugation by $abc$
sends the harmonic pair $(X_1,X_2)$ exactly to $(X_2,X_3)$.
Thus the preceding classification covers both harmonic pairs and
every source core.

Let $E$ be the simultaneous image of $N$ on all primitive
five-dimensional chain blocks of all thirty marked maps: the five
evaluation pairs and the six $\widehat B/F$ peripheral markings.
This marking enlargement is essential for retaining a source-normal
kernel. Because $N$ is $\widehat B$-normal, each marked map has
the same image on $N$ as its unmarked map, so the preceding simple
image classification applies to every projection.
Its kernel upon passing to the proper blocks is a finite three-group:
it is a subgroup of the group preserving the boundary/cycle/chain
filtrations and inducing the identity on all graded pieces.
Its proper image is a subdirect product of centreless nonabelian
simple groups, hence a direct product of diagonal simple factors
and in particular perfect. Therefore $E^{\rm ab}$ is a three-group,
and $E$ has no quotient of order two.

On the other hand $q_{12}(N)=\langle5I\rangle\simeq C_2$.
Goursat's lemma applied to the actual joint image
\[
                      N\longrightarrow E\times\langle5I\rangle
\]
now gives the full product: a proper subdirect image would furnish
a common nontrivial quotient, necessarily $C_2$, which $E$ lacks.
Consequently there is an actual source element $A\in N$ such that
\[
 q_{12}(A)=5I,\qquad
 A\text{ has identity image on every primitive full-chain block}.
 \tag{FP.56.1}
\]
This correction remains in $F''$ and in every original normal source
core containing $S$. It also retains every previous nilpotent-image
condition for $N$. If a literal finite expression is desired, it
can be chosen from the dense algebraic normal closure of $S$,
since (FP.56.1) imposes only a finite quotient condition.
No lift of an arbitrary finite $f$ assignment is being inferred.
The accepted R2 construction may therefore be applied with this
corrected source element; the obstruction in (FP.55) rules out the
old seed itself, but does not rule out the fourth-power target.

Complete classifications, Hermitian conjugacies, and the exact
source-core checks are respectively
\path{classify_d4_proper_seed_images.py},
\path{check_d4_proper_simple_conjugacies.py}, and
\path{check_d4_global_correction_inputs.py}, with their
JSON certificates and captured outputs. The reference simplicity
calculation is independently duplicated in the lifting audit's
\path{check_d4_proper_stabilizer.py}.

\paragraph{FP.57 (Proved: the reverse-prime full-chain target).}
Set $D_3=C''(C')^3$. Then $C'/D_3\simeq(C_3)^5$,
$[C:D_3]=1944$, and $D_3$ is free of rank $973$.
Over $\mathbb F_4=\mathbb F_2[\omega]$, $\omega^2+\omega+1=0$,
the same literal five-generator Fox construction gives an actual
$L$-representation on the $243$ five-dimensional chain blocks
indexed by $\chi\in(\mathbb Z/3)^5$. Source-to-target transport
is $\chi\mapsto\chi T_g^{-1}$, with the Fox matrix evaluated
at the target character. All permutations are retained.
All six inverse and boundary identities, the braid and full-twist
relations, and harmonic conjugacy by $abc$ are checked exactly.
The $X_i$ have orders $6,6,6,12,12$.

The four cusp words in (FP.54) give total boundary dimension
$323$ and proper dimension $650$, adding to $973$.
The harmonic character action has $27$ fixed characters and
$27$ orbits of size eight. Its fixed proper-three-dimensional
characters have indices $41,43,77,79,125,161,202,238$, in
base-three indexing. No commuting-permutation or blockwise-word
assumption is used.

For $w\in(C')'$, the trivial Fox derivative is zero, and the
nontrivial full chain blocks detect its class modulo $D_3'D_3^2$:
\[
 J_{\operatorname{Ad}(w)}(\chi)=I+m_\chi d_\chi,
 \qquad d_\chi m_\chi=0.
\]
This applies to the retained witness with pentagon in
$D_2'\subset(C')'$. It does not assert faithfulness on the
whole characteristic-two Ferrari quotient. The structural
and all-$2^n$ proof is TF.1--TF.3 in
\path{agent_literature/Ferrari_third_power_Fox.tex}.
Exact matrices and filtrations are in
\path{check_ferrari_d3_fox.py} and \path{check_d3_filtration.py},
with captured outputs.

\paragraph{FP.58 (Refuted: the original cyclic seed, or its cube, has nilpotent normal image in the reverse-prime target).}
The same actual $S$ from (FP.52) has identity values in all
eighteen nonharmonic core evaluations. Every harmonic core
value has order twelve and an order-two character permutation
moving $216$ characters. These are genuinely semilinear values.
For $x=X_1$, $y=X_2$, the elements $S^4$ and $S^9$ have
orders three and four, respectively, but
\[
                     [S^4,xS^9x^{-1}]\ne1.
 \tag{FP.58.1}
\]
This commutator fixes the character indices and is nonidentity
on $52$ chain blocks. Its first nontrivial block has index $10$,
character $(1,0,1,0,0)$; there the proper quotient is identity,
but the full chain shear is nonzero. In a finite nilpotent group
two-primary and three-primary elements commute, contradicting
(FP.58.1).

Cubing the seed, while retaining $q_{12}(S^3)=5I$, does not
repair the problem. Put
\[
                  u=S^3,\qquad v=(yx)u(yx)^{-1}.
 \tag{FP.58.2}
\]
Both have order four, but $uv$ has order twelve in the whole
target. At index $37$, character $(1,0,1,1,0)$, the full-chain
product has order six; both the cycle-space and proper
three-dimensional products have order three. The permutation
of $u$ exchanges this character with index $42$, and the
calculation retains that transport. The product $uv$ fixes all
character indices, so these block order tests are valid images
of the displayed source product. A product of two two-elements
in a finite nilpotent group is a two-element, giving the second
contradiction.

The thirty evaluations are \path{check_d3_controlled_seed.py}.
The short literal certificate \path{check_d3_seed_obstruction.py}
reconstructs $S$, verifies the primary orders and transports,
and saves the five-by-five matrices. These are obstructions to
two particular source seeds. The normal-image classification
and comparison with the signed frame coordinate remain open
at this checkpoint; no full-GT counterexample or universal
failure of R2 is asserted.
